% Generated by Sphinx.
\def\sphinxdocclass{report}
\documentclass[a5paper,10pt,polish]{book}
\usepackage[utf8]{inputenc}
\DeclareUnicodeCharacter{00A0}{\nobreakspace}
\usepackage{cmap}
\usepackage[T1]{fontenc}
\usepackage{babel}
\usepackage{times}
\usepackage[Sonny]{fncychap}
\usepackage{longtable}
\usepackage{sphinx}
\usepackage{multirow}
\usepackage{polski}
\usepackage[margin=0.75in]{geometry}
\pagenumbering{arabic}
\pagestyle{plain}

\title{Konspekty Rekolekcji Weekendowych}
\date{Wrzesień 2014}
\release{testowe}
\author{Diakonia Rekolekcji Weekendowych}
\newcommand{\sphinxlogo}{\includegraphics{foska.png}\par}
\renewcommand{\releasename}{Wydanie}
\makeindex

\makeatletter
\def\PYG@reset{\let\PYG@it=\relax \let\PYG@bf=\relax%
    \let\PYG@ul=\relax \let\PYG@tc=\relax%
    \let\PYG@bc=\relax \let\PYG@ff=\relax}
\def\PYG@tok#1{\csname PYG@tok@#1\endcsname}
\def\PYG@toks#1+{\ifx\relax#1\empty\else%
    \PYG@tok{#1}\expandafter\PYG@toks\fi}
\def\PYG@do#1{\PYG@bc{\PYG@tc{\PYG@ul{%
    \PYG@it{\PYG@bf{\PYG@ff{#1}}}}}}}
\def\PYG#1#2{\PYG@reset\PYG@toks#1+\relax+\PYG@do{#2}}

\expandafter\def\csname PYG@tok@gd\endcsname{\def\PYG@tc##1{\textcolor[rgb]{0.63,0.00,0.00}{##1}}}
\expandafter\def\csname PYG@tok@gu\endcsname{\let\PYG@bf=\textbf\def\PYG@tc##1{\textcolor[rgb]{0.50,0.00,0.50}{##1}}}
\expandafter\def\csname PYG@tok@gt\endcsname{\def\PYG@tc##1{\textcolor[rgb]{0.00,0.27,0.87}{##1}}}
\expandafter\def\csname PYG@tok@gs\endcsname{\let\PYG@bf=\textbf}
\expandafter\def\csname PYG@tok@gr\endcsname{\def\PYG@tc##1{\textcolor[rgb]{1.00,0.00,0.00}{##1}}}
\expandafter\def\csname PYG@tok@cm\endcsname{\let\PYG@it=\textit\def\PYG@tc##1{\textcolor[rgb]{0.25,0.50,0.56}{##1}}}
\expandafter\def\csname PYG@tok@vg\endcsname{\def\PYG@tc##1{\textcolor[rgb]{0.73,0.38,0.84}{##1}}}
\expandafter\def\csname PYG@tok@m\endcsname{\def\PYG@tc##1{\textcolor[rgb]{0.13,0.50,0.31}{##1}}}
\expandafter\def\csname PYG@tok@mh\endcsname{\def\PYG@tc##1{\textcolor[rgb]{0.13,0.50,0.31}{##1}}}
\expandafter\def\csname PYG@tok@cs\endcsname{\def\PYG@tc##1{\textcolor[rgb]{0.25,0.50,0.56}{##1}}\def\PYG@bc##1{\setlength{\fboxsep}{0pt}\colorbox[rgb]{1.00,0.94,0.94}{\strut ##1}}}
\expandafter\def\csname PYG@tok@ge\endcsname{\let\PYG@it=\textit}
\expandafter\def\csname PYG@tok@vc\endcsname{\def\PYG@tc##1{\textcolor[rgb]{0.73,0.38,0.84}{##1}}}
\expandafter\def\csname PYG@tok@il\endcsname{\def\PYG@tc##1{\textcolor[rgb]{0.13,0.50,0.31}{##1}}}
\expandafter\def\csname PYG@tok@go\endcsname{\def\PYG@tc##1{\textcolor[rgb]{0.20,0.20,0.20}{##1}}}
\expandafter\def\csname PYG@tok@cp\endcsname{\def\PYG@tc##1{\textcolor[rgb]{0.00,0.44,0.13}{##1}}}
\expandafter\def\csname PYG@tok@gi\endcsname{\def\PYG@tc##1{\textcolor[rgb]{0.00,0.63,0.00}{##1}}}
\expandafter\def\csname PYG@tok@gh\endcsname{\let\PYG@bf=\textbf\def\PYG@tc##1{\textcolor[rgb]{0.00,0.00,0.50}{##1}}}
\expandafter\def\csname PYG@tok@ni\endcsname{\let\PYG@bf=\textbf\def\PYG@tc##1{\textcolor[rgb]{0.84,0.33,0.22}{##1}}}
\expandafter\def\csname PYG@tok@nl\endcsname{\let\PYG@bf=\textbf\def\PYG@tc##1{\textcolor[rgb]{0.00,0.13,0.44}{##1}}}
\expandafter\def\csname PYG@tok@nn\endcsname{\let\PYG@bf=\textbf\def\PYG@tc##1{\textcolor[rgb]{0.05,0.52,0.71}{##1}}}
\expandafter\def\csname PYG@tok@no\endcsname{\def\PYG@tc##1{\textcolor[rgb]{0.38,0.68,0.84}{##1}}}
\expandafter\def\csname PYG@tok@na\endcsname{\def\PYG@tc##1{\textcolor[rgb]{0.25,0.44,0.63}{##1}}}
\expandafter\def\csname PYG@tok@nb\endcsname{\def\PYG@tc##1{\textcolor[rgb]{0.00,0.44,0.13}{##1}}}
\expandafter\def\csname PYG@tok@nc\endcsname{\let\PYG@bf=\textbf\def\PYG@tc##1{\textcolor[rgb]{0.05,0.52,0.71}{##1}}}
\expandafter\def\csname PYG@tok@nd\endcsname{\let\PYG@bf=\textbf\def\PYG@tc##1{\textcolor[rgb]{0.33,0.33,0.33}{##1}}}
\expandafter\def\csname PYG@tok@ne\endcsname{\def\PYG@tc##1{\textcolor[rgb]{0.00,0.44,0.13}{##1}}}
\expandafter\def\csname PYG@tok@nf\endcsname{\def\PYG@tc##1{\textcolor[rgb]{0.02,0.16,0.49}{##1}}}
\expandafter\def\csname PYG@tok@si\endcsname{\let\PYG@it=\textit\def\PYG@tc##1{\textcolor[rgb]{0.44,0.63,0.82}{##1}}}
\expandafter\def\csname PYG@tok@s2\endcsname{\def\PYG@tc##1{\textcolor[rgb]{0.25,0.44,0.63}{##1}}}
\expandafter\def\csname PYG@tok@vi\endcsname{\def\PYG@tc##1{\textcolor[rgb]{0.73,0.38,0.84}{##1}}}
\expandafter\def\csname PYG@tok@nt\endcsname{\let\PYG@bf=\textbf\def\PYG@tc##1{\textcolor[rgb]{0.02,0.16,0.45}{##1}}}
\expandafter\def\csname PYG@tok@nv\endcsname{\def\PYG@tc##1{\textcolor[rgb]{0.73,0.38,0.84}{##1}}}
\expandafter\def\csname PYG@tok@s1\endcsname{\def\PYG@tc##1{\textcolor[rgb]{0.25,0.44,0.63}{##1}}}
\expandafter\def\csname PYG@tok@gp\endcsname{\let\PYG@bf=\textbf\def\PYG@tc##1{\textcolor[rgb]{0.78,0.36,0.04}{##1}}}
\expandafter\def\csname PYG@tok@sh\endcsname{\def\PYG@tc##1{\textcolor[rgb]{0.25,0.44,0.63}{##1}}}
\expandafter\def\csname PYG@tok@ow\endcsname{\let\PYG@bf=\textbf\def\PYG@tc##1{\textcolor[rgb]{0.00,0.44,0.13}{##1}}}
\expandafter\def\csname PYG@tok@sx\endcsname{\def\PYG@tc##1{\textcolor[rgb]{0.78,0.36,0.04}{##1}}}
\expandafter\def\csname PYG@tok@bp\endcsname{\def\PYG@tc##1{\textcolor[rgb]{0.00,0.44,0.13}{##1}}}
\expandafter\def\csname PYG@tok@c1\endcsname{\let\PYG@it=\textit\def\PYG@tc##1{\textcolor[rgb]{0.25,0.50,0.56}{##1}}}
\expandafter\def\csname PYG@tok@kc\endcsname{\let\PYG@bf=\textbf\def\PYG@tc##1{\textcolor[rgb]{0.00,0.44,0.13}{##1}}}
\expandafter\def\csname PYG@tok@c\endcsname{\let\PYG@it=\textit\def\PYG@tc##1{\textcolor[rgb]{0.25,0.50,0.56}{##1}}}
\expandafter\def\csname PYG@tok@mf\endcsname{\def\PYG@tc##1{\textcolor[rgb]{0.13,0.50,0.31}{##1}}}
\expandafter\def\csname PYG@tok@err\endcsname{\def\PYG@bc##1{\setlength{\fboxsep}{0pt}\fcolorbox[rgb]{1.00,0.00,0.00}{1,1,1}{\strut ##1}}}
\expandafter\def\csname PYG@tok@kd\endcsname{\let\PYG@bf=\textbf\def\PYG@tc##1{\textcolor[rgb]{0.00,0.44,0.13}{##1}}}
\expandafter\def\csname PYG@tok@ss\endcsname{\def\PYG@tc##1{\textcolor[rgb]{0.32,0.47,0.09}{##1}}}
\expandafter\def\csname PYG@tok@sr\endcsname{\def\PYG@tc##1{\textcolor[rgb]{0.14,0.33,0.53}{##1}}}
\expandafter\def\csname PYG@tok@mo\endcsname{\def\PYG@tc##1{\textcolor[rgb]{0.13,0.50,0.31}{##1}}}
\expandafter\def\csname PYG@tok@mi\endcsname{\def\PYG@tc##1{\textcolor[rgb]{0.13,0.50,0.31}{##1}}}
\expandafter\def\csname PYG@tok@kn\endcsname{\let\PYG@bf=\textbf\def\PYG@tc##1{\textcolor[rgb]{0.00,0.44,0.13}{##1}}}
\expandafter\def\csname PYG@tok@o\endcsname{\def\PYG@tc##1{\textcolor[rgb]{0.40,0.40,0.40}{##1}}}
\expandafter\def\csname PYG@tok@kr\endcsname{\let\PYG@bf=\textbf\def\PYG@tc##1{\textcolor[rgb]{0.00,0.44,0.13}{##1}}}
\expandafter\def\csname PYG@tok@s\endcsname{\def\PYG@tc##1{\textcolor[rgb]{0.25,0.44,0.63}{##1}}}
\expandafter\def\csname PYG@tok@kp\endcsname{\def\PYG@tc##1{\textcolor[rgb]{0.00,0.44,0.13}{##1}}}
\expandafter\def\csname PYG@tok@w\endcsname{\def\PYG@tc##1{\textcolor[rgb]{0.73,0.73,0.73}{##1}}}
\expandafter\def\csname PYG@tok@kt\endcsname{\def\PYG@tc##1{\textcolor[rgb]{0.56,0.13,0.00}{##1}}}
\expandafter\def\csname PYG@tok@sc\endcsname{\def\PYG@tc##1{\textcolor[rgb]{0.25,0.44,0.63}{##1}}}
\expandafter\def\csname PYG@tok@sb\endcsname{\def\PYG@tc##1{\textcolor[rgb]{0.25,0.44,0.63}{##1}}}
\expandafter\def\csname PYG@tok@k\endcsname{\let\PYG@bf=\textbf\def\PYG@tc##1{\textcolor[rgb]{0.00,0.44,0.13}{##1}}}
\expandafter\def\csname PYG@tok@se\endcsname{\let\PYG@bf=\textbf\def\PYG@tc##1{\textcolor[rgb]{0.25,0.44,0.63}{##1}}}
\expandafter\def\csname PYG@tok@sd\endcsname{\let\PYG@it=\textit\def\PYG@tc##1{\textcolor[rgb]{0.25,0.44,0.63}{##1}}}

\def\PYGZbs{\char`\\}
\def\PYGZus{\char`\_}
\def\PYGZob{\char`\{}
\def\PYGZcb{\char`\}}
\def\PYGZca{\char`\^}
\def\PYGZam{\char`\&}
\def\PYGZlt{\char`\<}
\def\PYGZgt{\char`\>}
\def\PYGZsh{\char`\#}
\def\PYGZpc{\char`\%}
\def\PYGZdl{\char`\$}
\def\PYGZhy{\char`\-}
\def\PYGZsq{\char`\'}
\def\PYGZdq{\char`\"}
\def\PYGZti{\char`\~}
% for compatibility with earlier versions
\def\PYGZat{@}
\def\PYGZlb{[}
\def\PYGZrb{]}
\makeatother

\begin{document}
\shorthandoff{"}
\maketitle
\tableofcontents
\phantomsection\label{index::doc}



\part{Odpowiedzialność za wspólnotę}
\label{miliardowice-2005-jesien/index:odpowiedzialnosc-za-wspolnote}\label{miliardowice-2005-jesien/index:spis-tresci}\label{miliardowice-2005-jesien/index::doc}\paragraph{Podstawowe informacje o rekolekcjach}
\begin{itemize}
\item {} \begin{description}
\item[{Data:}] \leavevmode\begin{itemize}
\item {} 
10 listopada 2005 roku

\end{itemize}

\end{description}

\item {} \begin{description}
\item[{Miejsce pierwszych rekolekcji:}] \leavevmode\begin{itemize}
\item {} 
Miliardowice, archidiecezja katowicka

\end{itemize}

\end{description}

\item {} \begin{description}
\item[{Miejsce powstania konspektów:}] \leavevmode\begin{itemize}
\item {} 
Knurów, archidiecezja katowicka

\end{itemize}

\end{description}

\item {} \begin{description}
\item[{Grupa docelowa:}] \leavevmode\begin{itemize}
\item {} 
Oaza parafialna, od preoazy do formacji przed III st.

\end{itemize}

\end{description}

\item {} \begin{description}
\item[{Opis rekolekcji:}] \leavevmode\begin{itemize}
\item {} 
Rekolekcje parafialne organizowane przez lokalną diakonie mające na celu integrację i zwiększenie świadomości odpowiedzialności za swoją wspólnotę.

\end{itemize}

\end{description}

\end{itemize}


\chapter{Spotkanie 1 - Życie jako odpowiedź na powołanie do służby}
\label{miliardowice-2005-jesien/spotkanie1:spotkanie-1-zycie-jako-odpowiedz-na-powolanie-do-sluzby}\label{miliardowice-2005-jesien/spotkanie1::doc}

\section{Wprowadzenie (nie)tylko dla animatora}
\label{miliardowice-2005-jesien/spotkanie1:wprowadzenie-nie-tylko-dla-animatora}
Zacznijmy od sprawdzenia podstawowych funkcji - takich jak cytowanie:
\begin{quote}

Każde stworzenie mówi i jest powołane do wyrażenia jakiegoś szczególnego aspektu myśli Boga. Tam odnajduje swoje imię i tożsamość; potwierdza i czyni bezpieczną swoją wolność i oryginalność
\begin{quote}

--  Jan Paweł II - \emph{Dokument Końcowy Kongresu Powołań w Rzymie}
\end{quote}
\end{quote}

Na początek warto przytoczyć kilka fragmentów, które rozjaśnią nam nieco czym będziemy się zajmować.  Warto oczywiście wykorzystać je w trakcie spotkania. Powinny być wcześniej przemyślane  przez animatora, to baza dzięki której nabierzemy pewności w temacie.
\begin{itemize}
\item {} 
„Ten prosty fakt zaistnienia powinien nade wszystko przepełnić wszystkich zachwytem i niezmierną wdzięcznością dla Tego, który w całkowicie darmowy sposób powołał nas z niczego, wymawiając nasze imię. Zrozumienie zatem, że \textbf{życie jest darem}, nie powinno jedynie wzbudzić postawy uznania, lecz powoli powinno sugerować pierwszą wielką odpowiedź na zasadnicze pytanie o jego sens: życie człowieka jest arcydziełem stwórczej miłości Boga i jest samo w sobie wezwaniem  do miłości. \textbf{Otrzymany dar  ze swojej  natury  dąży  do stania  się dobrem dawanym}. Miłość jest pełnym sensem życia. Bóg tak umiłował człowieka, że oddał za niego Swoje życie i uczynił go zdolnym do życia i miłowania na sposób boski.” (Dokument końcowy Kongresu Powołań)

\item {} 
„Jak świętość jest dla wszystkich ochrzczonych w Chrystusie, tak też istnieje specyficzne powołanie każdego żyjącego; jak ta pierwsza jest zaszczepiona we Chrzcie świętym, tak to drugie jest związane z samym faktem istnienia. \textbf{Powołanie jest myślą opatrznościową Stwórcy względem każdego stworzenia, jest Jego ideą-projektem, jakby pewnym  marzeniem,  które  jest w sercu Boga, ponieważ los stworzenia leży Mu na sercu}. Bóg-Ojciec pragnie, by było ono różne i specyficzne dla każdego żyjącego. Istota ludzka jest w rzeczy- wistości «powołana» do życia i, gdy zaczyna żyć, niesie i odnajduje w sobie obraz Tego, który ją powołał. Powołanie jest boską propozycją do zrealizowania  się według tego obrazu i jest  jedyne,  jednostkowe, niepowtarzalne,  i dlatego też taki  obraz jest niewyczerpalny. Każde stworzenie mówi i jest powołane do wyrażenia jakiegoś  szczególnego aspektu myśli Boga. Tam odnajduje swoje imię i tożsamość; potwierdza i czyni bezpieczną swoją wolność i oryginalność.” (Dokument końcowy Kongresu Powołań)

\item {} 
\textbf{„Częścią  kultury powołaniowej  jest: zdolność snucia marzeń i wielkich pragnień; zdumienie, które pozwala docenić piękno i wybrać  je dla samej jego wewnętrznej wartości, ponieważ czyni życie pięknym  i prawdziwym;  altruizm, który  jest nie tylko solidarnością w potrzebie, lecz rodzi się z odkrycia godności kogokolwiek z sióstr i braci}. Kulturze przyjemności, która ryzykuje straceniem z pola widzenia i unicestwieniem poważnych pytań, przeciwstawia się kultura zdolna do znalezienia odwagi i upodobania w wielkich pytaniach człowieka, odnoszących  się do jego przyszłości. \textbf{W rzeczywistości są wielkimi  te pytania,  które czynią wielkimi  również małe  odpowiedzi}. Nic i nikt  nie może zdławić w człowieku pytania o sens i pragnienia prawdy.” (Dokument końcowy Kongresu Powołań w Rzymie)

\end{itemize}

\begin{DUlineblock}{0em}
\item[] \emph{Uwagi:}
\item[] Ks. dr Marek Dziewiecki od kwietnia 1999 jest krajowym duszpasterzem powołań. Dlatego często tutaj padają jego słowa gdyż zakładam, że jest to jedna z najbardziej kompetentnych osób piszących w języku polskim na ten temat.
\item[] \emph{Przypomnienie:}
\item[] \textbf{Spotkania Ruchu Światło-Życie  mają charakter dzieleniowy, nie wykładowy. Mamy sami siebie ubogacać. Wartością spotkania nie ma być treść konspektu tylko spotkanie z drugim człowiekiem w duchu zaufania oraz podzielenie się samym sobą.}
\end{DUlineblock}


\section{Modltiwa początkowa}
\label{miliardowice-2005-jesien/spotkanie1:modltiwa-poczatkowa}
Dowolna z uwzględnieniem szczególnie Ducha Świętego. Modlitwę należy dostosować do poziomu grupy


\section{Wprowadzenie do spotkania grupowego}
\label{miliardowice-2005-jesien/spotkanie1:wprowadzenie-do-spotkania-grupowego}
Być może w tym miejscu warto zapytać grupę jakie są ich oczekiwania od tego czasu rekolekcyjnego – spełniłoby to rolę integracyjną zamiast dzielenia się takimi informacjami o sobie, które nie zawsze są interesujące  z punktu widzenia rekolekcji.

Odczytanie Fragmentu, a następnie umieszczenie go w widocznym miejscu, aby „towarzyszył” nam w tym spotkaniu:
\begin{quote}

Stwarzając człowieka  na swój obraz i podobieństwo  oraz nieustannie podtrzymując go w istnieniu, Bóg wpisuje w człowieczeństwo mężczyzny i kobiety powołanie, a więc zdolność i odpowiedzialność za miłość i wspólnotę. Miłość jest zatem podstawowym i wrodzonym powołaniem każdej istoty ludzkiej
\begin{quote}

--  Jan Paweł II - \emph{Adhortacja „Familiaris consortio”}
\end{quote}
\end{quote}

Fragment ten jest równocześnie wprowadzeniem do tematyki spotkania – zapytać konkretnie kogoś jak reaguje gdy widzi taki tekst, czy są to dla niego „frazesy” czy może nieżyciowe rozważania, a może czymś co poruszyło?


\section{Dar  życia: co z z nim zrobić? Czym  jest?}
\label{miliardowice-2005-jesien/spotkanie1:dar-zycia-co-z-z-nim-zrobic-czym-jest}
Bez odpowiedzi na pytania: czym jest moje życie, o czym świadczy to, że istnieję, po co żyję itp. nie ma możliwości pójścia dalej. Człowiek zazwyczaj albo unika tych pytań, bo są dla niego zbyt trudne, albo spycha w „przestrzeń oczywistości”.
\begin{itemize}
\item {} 
Jak wygląda nasza świadomość tego absolutnie nietrywialnego faktu, że istniejemy, czujemy, oddychamy, tworzymy? Czy zastanawialiśmy się kiedykolwiek nad tym?

\end{itemize}

Nasze życie umieszczone jest w czasie, pobyt tutaj na ziemi będzie trwał konkretną, skończoną, liczbę sekund. Nie daliśmy sobie życia sami, niewątpliwie otrzymaliśmy je od kogoś. Jest zatem darem. To zobowiązuje, aby nie zmarnować tak drogocennej rzeczy którą nam powierzono (na dodatek za darmo). Nasz czas ziemski jest zatem bardzo cenny – służy nam do pielęgnowania życia.
\begin{itemize}
\item {} 
Czy z świadomości daru życia wynika dla nas troska o nasze działania? Czy zadajemy sobie pytania w stylu „w jaki sposób wykorzystuje własny czas?” Czy szukamy sposobu aby jak najlepiej go wykorzystać?

\end{itemize}

To bardzo poważna i fundamentalna sprawa – wyrwanie się z myślenia pragmatycznego do myślenia o sensie własnego życia.

Realizować życie można na tyle sposobów ilu ludzi jest na świecie. Można próbować samemu absolutnie od podstaw poszukiwać jego sensu. To droga jednak raczej dla niewielu. Kościół głosząc Dobrą Nowinę o Chrystusie pokazuje bardzo wyraźny kierunek-propozycję dla nas.


\section{Powołanie:  co zrobić z moim życiem? Jaka jest moja droga?}
\label{miliardowice-2005-jesien/spotkanie1:powolanie-co-zrobic-z-moim-zyciem-jaka-jest-moja-droga}
Bóg wybiera każdego z nas. On nas stworzył i nie dokonał tego „ot tak sobie bez pomysłu”, dając nam dar życia miał konkretny zamysł. On chce byśmy byli piękni, abyśmy byli na Jego podobieństwo.
\begin{quote}

Miłość Boga Ojca, która stworzyła każdego człowieka, sprawia, że  nikt nie powinien czuć się niepotrzebny, samotny czy zapomniany, gdyż jest powołany do życia w miłości  według  planu Bożego, ułożonego właśnie dla niego, czyli zgodnie  z jego niepowtarzalną historią, aspiracjami oraz możliwościami. Poszczególny człowiek ma szansę czuć się kochany, bezpieczny, zrealizowany i w pełni usatysfakcjonowany tylko wtedy, gdy odkryje i przyjmie Boży zamysł, Boże marzenie na temat jego życia i przyszłości. Z tego powodu to właśnie Bóg Ojciec jest najlepszym wychowawcą człowieka,  gdyż tylko On może w pełni wydobyć z serca człowieka to, co On sam tam umieścił. Tylko Bóg może skutecznie pomóc człowiekowi, by stał się on najpiękniejszą wersją samego siebie.
\begin{quote}

--  Ks. Marek Dziewiecki
\end{quote}
\end{quote}
\begin{itemize}
\item {} 
Czy zgadzamy  się z tymi słowami? Czy możemy  się podzielić doświadczeniem właśnie takiego działania Boga, że On wydobył z nas rzeczy które nas samych wprowadziły w zdumienie?

\item {} 
Jak wygląda nasza wiara w to, że ta „najpiękniejsza wersja samego siebie” istnieje i mamy obowiązek  się o nią starać?

\end{itemize}

Powstaje perspektywa, że ten dar życia o którym mówiliśmy na początku ma konkretny cel. Nie jesteśmy tym czym moglibyśmy być. Jest ktoś prócz nas samych komu zależy na tym abyśmy tym byli. Ma on możliwości aby nam pomóc, musimy jednak z nim współpracować.


\section{Służba: postawa do naśladowania}
\label{miliardowice-2005-jesien/spotkanie1:sluzba-postawa-do-nasladowania}
Bóg przyszedł na ziemie jako człowiek, aby nam pomóc. Ukazał nam doskonałość do której mamy dążyć swoim życiem. Mówił wielokrotnie, że mamy go naśladować.
\begin{itemize}
\item {} \begin{description}
\item[{Przeczytać Flp 2,5-11}] \leavevmode\begin{itemize}
\item {} 
Ogołocił, przyjął postać sługi, uniżył samego siebie, był posłuszny

\end{itemize}

\end{description}

\item {} \begin{description}
\item[{Przeczytać J 13,1-17}] \leavevmode\begin{itemize}
\item {} 
Obmywał nogi

\end{itemize}

\end{description}

\item {} 
Dlaczego apostołowie mimo, że tyle miesięcy wędrowali z Jezusem i słuchali  jego nauk nie rozumieli tego  co uczynił? Czy i dlaczego postawa służby jest dla nas tak „nieintuicyjna”?

\end{itemize}

Służba jest wyrazem miłości. Ukierunkowuje nas na drugiego, a nie na siebie, wyklucza egoizm który jest przecież w nas „wrodzony”.
\begin{itemize}
\item {} 
Jakie mamy największe problemy w służbie?

\end{itemize}

Aby służyć innym (obdarowywać ich swoją miłością) musimy umieć także otrzymywać – inaczej nie damy możliwości naszym bliźnim posługi. To bardzo ważna kwestia.
\begin{itemize}
\item {} 
Czy doświadczyliśmy w naszym życiu czegoś takiego, że z powodu nieumiejętności (swojej/cudzej) otrzymywania powstało zniechęcenie u osoby dającej? (np. z powodu fałszywej pokory)

\end{itemize}


\section{Powołanie  do służby konsekwencją daru życia}
\label{miliardowice-2005-jesien/spotkanie1:powolanie-do-sluzby-konsekwencja-daru-zycia}\begin{quote}

Zrealizować we własnym życiu Boże powołanie to stać się wiernym uczniem Chrystusa, to stać się przedłużeniem Jego słów i Jego czynów, to stać się przedłużeniem Jego obecności na tej ziemi. To doprowadzić do sytuacji, w której za św. Pawłem możemy powiedzieć, że już nie ja żyję, lecz żyje we mnie Chrystus. To oddać Jezusowi do całkowitej dyspozycji samego siebie, aby poprzez nas Jego serce i Jego ręce mogły nadal obejmować małych i zagubionych, chorych i odrzuconych przez ludzi, słabych i grzeszników.
\begin{quote}

--  Ks. Marek Dziewiecki
\end{quote}
\end{quote}

Oddanie Jezusowi samego siebie do całkowitej dyspozycji  to nic innego jak stanie się Jego narzędziem, a sługą dla bliźniego.

Nie jednak tak, że od razu będziemy umieli do \emph{„całkowitej dyspozycji”}. Przez całe nasze życie walczymy o to aby tak  było. Decydując się zatem na służbę drugiemu musimy zdawać sobie sprawę z tego, że nie będziemy potrafili  mu służyć doskonale! Będziemy ranić własny egoizmem, własną grzecznością, będziemy rozczarowywać!

Kryje się tutaj niebezpieczna pułapka rezygnacji z służby dla pozornie większego dobra bliźniego!
\begin{itemize}
\item {} 
Czy zauważamy w naszym życiu zniechęcenie spowodowane świadomością naszej niedoskonałości, tym, że decydując  się na miłość automatycznie wprowadzamy ból? Jak z tym walczymy?

\end{itemize}

To właśnie jest Kościół – grzesznik grzesznikowi służy. Zamiast oddalać tą prawdę  jak najdalej  od osób które  kochamy powinniśmy razem z nimi zgłębiać tą tajemnice, mimo, że często  jest  to bolesne. Wtedy wyrażamy swoją największą troskę o tego człowieka, a nie wtedy gdy usuwamy się i pozostawiamy go samotnym „aby nie cierpiał”.
\begin{itemize}
\item {} 
Czy przeżywamy swoje powołanie jako radość z możliwości służenia, czyli upodabniania  się do Jezusa? Czy widzimy, że miłośc najpełniej wyraża się w służbie, czyli czynie?

\end{itemize}


\section{Zastosowanie}
\label{miliardowice-2005-jesien/spotkanie1:zastosowanie}
Metoda Ruchu Światło-Życie nie pozwala na to, abyśmy pozostawali tylko na poziomie wiedzy  i informacji. Każde spotkanie  powinno zakończyć się
„zastosowaniem”, czyli konkretnym wprowadzeniem w życie treści.

Propozycja, aby każdy z nas  podjął  teraz jedno małe postanowienie w związku z omawianym tematem, które liczy, że będzie  mógł zrealizować jeszcze na tych rekolekcjach.


\section{Modlitwa końcowa}
\label{miliardowice-2005-jesien/spotkanie1:modlitwa-koncowa}
Dowolna w intencji wytrwałości w naszych postanowieniach.


\part{Duchowość człowieka}
\label{babice2006-wiosna-knurow/index::doc}\label{babice2006-wiosna-knurow/index:duchowosc-czlowieka}\paragraph{Podstawowe informacje o rekolekcjach}
\begin{itemize}
\item {} \begin{description}
\item[{Data:}] \leavevmode\begin{itemize}
\item {} 
24 marca 2006 roku

\end{itemize}

\end{description}

\item {} \begin{description}
\item[{Miejsce pierwszych rekolekcji:}] \leavevmode\begin{itemize}
\item {} 
Babice, diecezja gliwicka

\end{itemize}

\end{description}

\item {} \begin{description}
\item[{Miejsce powstania konspektów:}] \leavevmode\begin{itemize}
\item {} 
Knurów, archidiecezja katowicka

\end{itemize}

\end{description}

\item {} \begin{description}
\item[{Grupa docelowa:}] \leavevmode\begin{itemize}
\item {} 
Oaza parafialna, od preoazy do formacji przed III st.

\end{itemize}

\end{description}

\item {} \begin{description}
\item[{Opis rekolekcji:}] \leavevmode\begin{itemize}
\item {} 
Rekolekcje parafialne organizowane przez lokalną diakonie mające na celu integrację i zwiększenie świadomości odpowiedzialności za swoją wspólnotę.

\end{itemize}

\end{description}

\end{itemize}


\chapter{Spotkanie 1 - Duchowa kondycja człowieka}
\label{babice2006-wiosna-knurow/spotkanie1:spotkanie-1-duchowa-kondycja-czlowieka}\label{babice2006-wiosna-knurow/spotkanie1::doc}

\section{Wprowadzenie dla Animatora}
\label{babice2006-wiosna-knurow/spotkanie1:wprowadzenie-dla-animatora}
Według Instrumentum laboris VIII Zwyczajnego Synodu Biskupów z 1990r. słowo ”duchowość oznacza pewien «styl życia», albo «życie według Ducha» (por. Rz 8,9), «postępowanie według Ducha» (por. Rz 8,4) {[}...{]}, zespół po- staw serc  pobudzonych przez Ducha, a mianowicie  przekonań, motywacji, decyzji” (nr 21). \textbf{Duchowość więc jest zbiorem postaw, przez które, zgodnie z myślą ks. kard. Karola  Wojtyły, rozumie się zajęcie sta- nowiska wobec kogoś lub  czegoś, nie będące jeszcze działaniem, a będące już czymś więcej niż świadomością i zarazem gotowością do zadziałania  według zajętego stanowiska.}

Spróbujmy „wyizolować” siebie i spojrzeć na nas w samotności (patrz Namiot Spotkania). Pokazać,  że samotność  jest nieodłączną częścią życia. Każdy nasz wybór dokonuje się w nagiej samotności.....spróbujmy dotknąć tego doświadczenia, tego jak się z nim czujemy.


\section{Wstęp. Dzielenie  się owocami dnia}
\label{babice2006-wiosna-knurow/spotkanie1:wstep-dzielenie-sie-owocami-dnia}\begin{itemize}
\item {} 
Czy ciężko  było Ci rozmawiać przez ten długi  czas  z drugim człowiekiem? Co przychodziło łatwo, a co było trudniejsze?

\item {} 
Co odkryłeś/odkryłaś w sobie dzięki temu czasowi?

\end{itemize}

\emph{Nie śpieszyć się z tym  punktem! Niech ludzie otwierają  się na siebie. Niech od razu będzie tutaj zaznaczone, że bogactwem  tego spotkania jest to czym sami się dzielimy.}


\section{Czym  jest duchowość?}
\label{babice2006-wiosna-knurow/spotkanie1:czym-jest-duchowosc}
Mamy zamiar rozmawiać o kondycji człowieka. Na osobę człowieka składa się wiele spraw, moglibyśmy wyróżnić: ciało, psychike, emocje i duchowość. Każdy z nas może własnym doświadczeniem potwierdzić, że mamy te rzeczy w sobie i są one nasze własne.  Dlaczego jednak właśnie duchowości chcemy poświęcić tyle czasu?

(Podajemy całej grupie stopniowo poniższe tezy. Próbujemy wspólnie je wytłumaczyć, milło byłoby gdyby było to oparte na własnym doświadczeniu. Jest to zasadnicza  część spotkania.) \footnote{
Należy mieć przygotowane wydrukowane definicje (do pobrania w wersji html)
}
\begin{itemize}
\item {} \begin{description}
\item[{Ciało}] \leavevmode\begin{itemize}
\item {} 
„Jeśli ktoś kontaktuje się z samym sobą jedynie poprzez ciało, to ono nie może mu wyjaśnić tajemnicy  człowieka, gdyż ciało nie wie, kim jest człowiek i po co on żyje. Ciało odczuwa jedynie własne potrzeby i popędy. Pozostawione samemu sobie ciało staje się rodzajem  nowotworu, który podporządkowuje  sobie ca- łego człowieka i żyje jego kosztem. Człowiek jedynie cielesny nie może więc zrozumieć samego siebie i nie może odpowiedzialnie kierować swoim życie. Człowiek cielesny to człowiek podlegający dyktaturze ciała.”

\end{itemize}

\end{description}

\item {} \begin{description}
\item[{Psychika}] \leavevmode\begin{itemize}
\item {} 
„Człowiek nie może zrozumieć samego siebie jeśli kontaktuje się jedynie z własną psychiką, a więc z własnymi strategiami myślenia czy własnymi sposobami przeżywania emocjonalnego. Umysł ludzki posiada zdolność myślenia ale nie jest źródłem wiedzy o tajemnicy człowieka. Umysł może tę tajemnicę odkryć ale nie może jej własną mocą ustalić czy wymyślić. Pozostawiony własnej spontaniczności umysł ludzki nie szuka prawdy lecz wygod- nych dla danego człowieka półprawd a nawet iluzji.  Człowiek, dla którego jego własne myślenie jest najwyższym i jedynym au- torytetem, kieruje się tzw. logiką prywatną po to, by nie odkryć całej o prawdy o swoim życiu, zwłaszcza tej trudnej i stawiającej jasne wymagania. Ale kierując się jedynie subiektywnym myśleniem człowiek wyrządza sobie zwykle wielką krzywdę, oszukuje samego siebie i pozbawia  się szansy na rozwój.”

\end{itemize}

\end{description}

\item {} \begin{description}
\item[{Emocje}] \leavevmode\begin{itemize}
\item {} 
„Także z perspektywy emocji człowiek nie może odkryć kim jest i po co żyje. Człowiek jedynie emocjonalny nie wie, jaką postawę powinien on zająć wobec własnej rzeczywistości a nawet wobec samych emocji.  Emocje  pozostawione własnej  spontaniczności kierują się własną logiką: szukają jedynie dobrego samopoczucia. Natychmiast i za wszelką cenę. Nawet za cenę niszczenia człowieka, którego są częścią. Wszelkie uzależnienia od osób, rzeczy czy substancji chemicznych są bezpośrednią konsekwencją dyktatury emocji, na którą skazuje się człowiek jedynie emocjonalny.”

\end{itemize}

\end{description}

\item {} \begin{description}
\item[{Duchowość}] \leavevmode\begin{itemize}
\item {} 
„Niezwykłością sfery duchowej jest fakt, że w tej sferze i tylko w tej  sferze człowiek może zapytać się o własną tajemnicę nie z perspektywy cząstkowej (cielesnej czy psychicznej)  lecz z perspektywy całego człowieka. \textbf{I tylko w tej sferze może znaleźć ostateczną odpowiedź, która stanie się podstawą zajęcia dojrzałej  postawy wobec życia.} Duchowość zaczyna się zatem dopiero wtedy, gdy człowiek wznosi się \textbf{ponad swoje ciało i ponad swoją psychikę, aby  postawić sobie pytanie  o całą swoją  rzeczywistość}:  o to, kim jest  i po co żyje, o to, jaki sens ma jego ciało, jego myślenie i jego emocje i on sam.

\end{itemize}

\end{description}

\end{itemize}

Nie  ma więc wolności bez duchowości, gdyż człowiek nie może odpowiedzialnie kierować sobą i swoim życiem dopóki nie wie, kim jest i po co żyje. \textbf{Człowiek  bez duchowości to człowiek z definicji uzależniony od jednej ze swoich sfer cząstkowych albo od nacisków płynących ze środowiska zewnętrznego}. To człowiek, który żyje jak małe dziecko, czyli nieświadomie i jedynie w oparciu o bodźce, które oddziałują na niego w danym momencie.”


\section{Pytanie  o naszą kondycję duchową}
\label{babice2006-wiosna-knurow/spotkanie1:pytanie-o-nasza-kondycje-duchowa}
Na ile potrafimy spoglądać na siebie z perspektywy duchowej? To znaczy: Patrzeć na całą naszą osobę - nie rozdzielać nas na kilka, nie mających ze sobą nic wspolnego, rzeczywistości?

Wypisując na kartce odpowiedzi pytamy:
\begin{itemize}
\item {} 
W jakich momentach życia takie podejście przychodzi nam najłatwiej?

\end{itemize}

Patrząc na uzyskaną kartkę próbujemy znaleźć jakieś prawidłowości.
\begin{itemize}
\item {} 
Jak myślimy: dlaczego akurat w takich momentach? co mają one wspólnego?

\end{itemize}


\section{Co możemy z własną duchowością zrobić?}
\label{babice2006-wiosna-knurow/spotkanie1:co-mozemy-z-wlasna-duchowoscia-zrobic}
Postawa duchowa jest „najważniejsza” wtedy gdy próbujemy odpowiedzieź sobie na te najistotniejsze pytania: „kim jestem? jaki jestem? po co jestem?” Co zrobić, aby w tym momencie potrafić patrzeć z tej perspektywy? Warto zwrócić uwagę jakie to są  momenty  życia? (zazwyczaj  momenty  sporego smutku i „rozchwianego stanu emocjonalnego”)
\begin{itemize}
\item {} 
Jakich „rad” udzielilibyśmy sami sobie teraz na tym spotkaniu, aby na takie pytania umieć odpowiadać dojrzale?

\end{itemize}

Podpowiedzi (kierować grupę na to):
\begin{enumerate}
\item {} 
rozwiazanie „najprostsze”: pytania egzystencjalne stawiać sobie wtedy gdy jest dobrze...gdy potrafimy sprawnie myśleć, gdy mamy ustabilizowane emocje, gdy nic „ nie leży nam” na psychice. Takie pytania wtedy nie przychodzą same, należy je „prowokować”

\item {} 
rozwiązanie wymagające pozornie więcej: na każdą sprawe  własnego życia spoglądać przez pryzmat duchowości, wtedy gdy przyjdą te pytania najważniejsze „mimochodem” podejdziemy do nich tak samo jak do każdych wcześniej

\end{enumerate}

\begin{DUlineblock}{0em}
\item[] (Najlepiej oczywiście połączyć oba te rozwiązania.)
\item[] Co jest  potrzebne, aby realizować pierwszą możliwość? Każdy ma zapewne swoje doświadczenie,  że sami z siebie „gdy jest dobrze” takich pytań sobie nie stawiamy.
\end{DUlineblock}
\begin{itemize}
\item {} 
Czy mamy kogoś kto „czuwa” nad tym, abyśmy sobie takie pytania zadawali? Kto  to jest? Czy czujemy wdzięczność za jego starania?

\end{itemize}

Współczesny świat często próbuje nam powiedzieć,  że pytanie o sprawy najważniejsze drugiego człowieka jest naruszaniem jego prywatności, czymś złym. Nie wypada pytać o coś osobistego,  o coś co dotyczy wnętrza.

Jednak osoba która nad nami czuwa nie robi tego 24h/dobę. Są chwile kiedy jesteśmy  tylko  z własną samotnością (mimo, że koło nas może być tłum).  Te chwile mogą nam wiele powiedzieć o stanie naszej kondycji du- chowej. Chodzi o samotność która istnieje między ludźmi (Jezus zawsze jest z nami, ale o tym dopiero jutro).
\begin{itemize}
\item {} 
Czy odczuwamy taką samotność? Jeżeli możemy to podajmy jakiś jej przykład.

\end{itemize}

Nieraz tak jest. . . np. idziemy na ważną dla nas rozmowę do najlepszego przyjaciela, mówimy mu otwarcie wiele rzeczy, gdy on odpowiada: „rozumiem Cię” możemy  się tylko gorzko uśmiechnąć. Gdzieś czujemy, \textbf{że on rozumieć do końca nie może, że to jest nasze osobiste i cokolwiek zrobimy takim  pozostanie}. Jest coś z czym musimy zostać samotni.
\begin{itemize}
\item {} 
Czy zdażyło nam się w czasie swojej samotności samemu wyjść  do drugiego człowieka,  aby  z nim  porozmawiać o rzeczach ważnych? W  jaki  sposób zrodziła  się w nas ta  chęć? Czy było to takie  „samo z siebie” czy jednak  pod wpływem impulsu?

\item {} 
Jak wykorzystujemy własny czas samotności?  Czy przeżywamy go jako czas przygotowawczy  do powrotu do ludzi?

\end{itemize}

Zauważmy,  że odpowiedź na powyższe pytania świadczą właśnie o naszej kondycji duchowej. Jeżeli staramy się patrzeć przez pryzmat duchowości to dlatego, że ktoś nas do tego „sprowokował”, a więc doświadczyliśmy, że te prowokowanie jest dobre i chcemy sami  innych także  zachęcać do tego - odpowiedzialność.

Istotą  tego o czym teraz mówimy jest  to czy potrafimy wychodzić do drugiego człowieka sami z siebie. . . to szalenie ważne, aby mieć taką umiejętność, dopóki ludzie będą ją posiadali i będą sobie ją przekazywać kondycja duchowa człowieka będzie „stabilna”.
\begin{itemize}
\item {} 
A  co jest potrzebne,  aby realizować drugą możliwość? By na każdą chwile życia spoglądać przez pryzmat duchowości?

\end{itemize}

Przede wszystkim świadomość wagi jaką ma duchowość (patrz. początek spotkania!). Ona natomiast powinna nas motywać do pracy nad sobą, do przypominania sobie nieustannie, pernamentnie kim jest człowiek.


\section{Mądra odwaga stawiania sobie pytań}
\label{babice2006-wiosna-knurow/spotkanie1:madra-odwaga-stawiania-sobie-pytan}
Chodzi o to, aby wspólnie dojść do tego, że jeżeli będziemy trzymać się zasad o których mówiliśmy wyżej znika lęk stawiania sobie pytań trudnych. Nie są one już zagrożeniem, a czymś co rozwija naszą osobowość. Dla człowieka który duchowość swoją traktuje „na zawołanie” niektóre pytania przerażają - szczególnie te na które odpowiedź należy szukać dłużej niż „przez chwilę”.
\begin{itemize}
\item {} 
Czy mamy odwagę stawiać  sami sobie egzystencjalne  pytania?

\item {} 
Czy widzimy taką potrzebę? Co jest takiego niebezpiecznego w tych pytaniach?

\end{itemize}

Niebezpieczne jest  odkrycie tego kim jesteśmy  naprawdę,  bez  masek. Pytania najważniejsze, przed którymi  stajemy samotnie,  muszą odkrywać prawdę o nas. Nie jest to niczym przyjemnym. Czasami szukującym.  Jest to „utrudnianie” sobie życia - choć tak naprawdę to dopiero jest prawdziwą egzystencją - bo świadomą.
\begin{itemize}
\item {} 
Co robimy z tą prawdą o sobie? Czy zawsze mamy świadomość, że nigdy nie poznaliśmy całej prawdy? Czy próbujemy nieustannie się siebie pytać o to kim jesteśmy?

\end{itemize}

Jedno z najważniejszych wniosków - niepozwolić by odpowiedzi na pytania o nas samych stały się naszymi  „prywatnymi dogmatami”, czymś do czego doszliśmy i już nie chcemy tego zmieniać - abyśmy nie wyeliminowali z własnego życia rozwoju! Tajemnice człowieka trzeba nam zgłębiać przez całe życie.

Na zakończenie przeczytajmy:
\begin{quote}

Człowiek  zmienia się od kołyski  aż po grób. Kiedy więc jest sobą?
\begin{quote}

--  Stefan Kisielewski
\end{quote}
\end{quote}


\section{Zastosowanie}
\label{babice2006-wiosna-knurow/spotkanie1:zastosowanie}
Metoda Ruchu Światło-Życie nie pozwala na to, abyśmy pozostawali tylko na poziomie wiedzy  i informacji. Każde spotkanie  powinno zakończyć się
„zastosowaniem”, czyli konkretnym wprowadzeniem w życie treści.

Propozycja, aby każdy z nas  podjął  teraz jedno małe postanowienie w związku z omawianym tematem, które liczy, że będzie  mógł zrealizować jeszcze na tych rekolekcjach.


\section{Modlitwa końcowa}
\label{babice2006-wiosna-knurow/spotkanie1:modlitwa-koncowa}
Dowolna w intencji wytrwałości w naszych postanowieniach  i za te rekolekcje.


\chapter{Spotkanie 2 - Boże spojrzenie na człowieka}
\label{babice2006-wiosna-knurow/spotkanie2:spotkanie-2-boze-spojrzenie-na-czlowieka}\label{babice2006-wiosna-knurow/spotkanie2::doc}

\section{Wprowadzenie dla Animatora}
\label{babice2006-wiosna-knurow/spotkanie2:wprowadzenie-dla-animatora}\begin{quote}

Człowiek nie może żyć bez miłości. Człowiek pozostaje dla siebie istotą  niezrozumiałą,  jego życie jest pozbawione  sensu, jeśli nie objawi mu się Miłość, jeśli nie spotka się z Miłością, jeśli jej nie dotknie i nie  uczyni  w i jakiś  sposób swoją,  jeśli  nie  znajdzie w i niej żywego uczestnictwa.  I dlatego właśnie Chrystus - Odkupiciel,  (...)  objawia w i pełni  człowieka  człowiekowi.  To jest ów - jeśli tak wolno się wyrazić - ludzki wymiar Tajemnicy Odkupienia.  Człowiek  odnajduje w i nim swoją właściwą wielkość, godność i wartość swego człowieczeństwa.  Człowiek zostaje w i Tajemnicy Odkupienia na nowo potwierdzony, niejako wypowiedziany na nowo. Stworzony na nowo! Nie ma już Żyda ani poganina, nie ma już niewolnika ani człowieka wolnego, nie ma już mężczyzny ani kobiety, wszyscy bowiem jesteście kimś jednym w i Chrystusie Jezusie. Człowiek, który chce zrozumieć siebie do końca - nie wedle jakiś tylko doraźnych, częściowych, czasem powierzchownych, a nawet pozornych kryteriów i miar swojej własnej istoty  - musi ze swoim niepokojem, niepewnością, a także słabością i grzesznością,  ze swoim życiem i śmiercią, przybliżyć się do Chrystusa. Musi niejako w i Niego wejść z sobą  samym, musi sobie przyswoić, zasymilować całą rzeczywistość Wcielenia i Odkupienia,  aby siebie odnaleźć. Jeśli dokona się w i człowieku ów dogłębny proces, wówczas owocuje on nie tylko uwielbieniem Boga, ale także głębokim  zdumieniem  nad sobą samym. Jakąż wartość musi mieć w i oczach Stwórcy człowiek,  skoro zasłużył na takiego i tak potężnego Odkupiciela, skoro  Bóg Syna swego Jednorodzonego dał, ażeby on, człowiek, nie zginął, ale miał życie wieczne.

Właśnie owo głębokie zdumienie  wobec wartości i godności człowieka nazywa się Ewangelią, czyli Dobrą Nowiną.
\begin{quote}

--  Jan Paweł II - \emph{Redemptor hominis}
\end{quote}
\end{quote}

Ten tekst oddaje wszystko o czym chcemy mówić tego dnia.


\section{Wstęp}
\label{babice2006-wiosna-knurow/spotkanie2:wstep}
Wczoraj  cały dzień mówiliśmy o kondycji duchowej  człowieka. Próbowaliśmy dokładnie przyjrzeć się  samym sobą i zastanowić się kim jesteśmy. Wieczorem zaczeliśmy mówić już o Bogu który w i te nasze życie wkracza. Chcemy teraz spojrzeć jeszcze raz na nas samych, ale zadając sobie pytanie: co „obecność Boga” zmienia w i wczorajszych rozważaniach.

\emph{Czytamy tekst z wprowadzenia - fragment encykliki Redemptor Hominis.}

Będziemy odczytywać teraz ten piękny tekst zdanie po zdaniu odnosząc do naszego tematu.


\section{Jezus nadaje sens ludzkiemu życiu}
\label{babice2006-wiosna-knurow/spotkanie2:jezus-nadaje-sens-ludzkiemu-zyciu}\begin{quote}

Człowiek nie może żyć bez miłości. Człowiek pozostaje dla siebie istotą niezrozumiałą, jego życie jest pozbawione sensu, jeśli nie objawi mu się Miłość, jeśli  nie spotka się  z Miłością, jeśli  jej nie dotknie i nie uczyni w i jakiś sposób swoją, jeśli nie znajdzie w i niej żywego uczestnictwa.
\end{quote}
\begin{itemize}
\item {} 
O jaką miłość tutaj chodzi? Jak odczytujemy to, że zrozumieć siebie możemy tylko przez spotkanie z Miłością?

\item {} 
Czy my w i miłości potrafimy odnaleźć pełnię swojego człowieczeństwa?  Czy taką „range” jej dajemy?

\end{itemize}

To trochę tak jakby chcieć „poznawać”  mechanizm np. drukarki  bez próby drukowania na niej czegokolwiek. Człowiek jest stworzony do miłości, dlatego tylko poprzez kontakt z nią (czyli odpowiadanie miłością) może sam siebie odkrywać.


\section{Objawienie  w i pełni człowieka człowiekowi}
\label{babice2006-wiosna-knurow/spotkanie2:objawienie-w-i-pelni-czlowieka-czlowiekowi}\begin{quote}

I dlatego właśnie Chrystus-Odkupiciel, (...)  objawia w i pełni człowieka człowiekowi. To jest ów - jeśli  tak wolno się wyrazić - ludzki wymiar Tajemnicy Odkupienia. Człowiek odnajduje w i nim swoją właściwą wielkość, godność i wartość swego człowieczeństwa. Człowiek zostaje w i Tajemnicy Odkupienia na nowo potwierdzony, niejako wypowiedziany na nowo. Stworzony na nowo! Nie ma już Żyda ani poganina, nie ma już niewolnika ani człowieka wolnego, nie  ma już  mężczyzny ani kobiety, wszyscy bowiem jesteście kimś jednym w i Chrystusie Jezusie.
\end{quote}
\begin{itemize}
\item {} 
Czy czytając powyższe zdania potrafimy tą wielkość,  godność odnaleźć w i obrazie „Ecce Homo”?

\item {} 
Na czym ona polega, jak my ją odczuwamy?

\item {} 
Czy i my potrafimy przeżywać podobną radość myśląc o tym, że Tajemnica Odkupienia „stwarza nas na nowo”?

\end{itemize}

Zwróćmy uwagę na to, że Bóg przychodzi  po to,a by pokazać człowiekowi kim on sam jest! To objawienie w i pełni człowieka człowiekowi oznacza, że zagubiliśmy „wcześniej”  własną tożsamość - nie wiemy kim jesteśmy. Dla Boga było zatem bardzo ważne,a byśmy się tego dowiedzieli.


\section{Odkrycie  dogłębnej prawdy  o sobie}
\label{babice2006-wiosna-knurow/spotkanie2:odkrycie-doglebnej-prawdy-o-sobie}\begin{quote}

Człowiek, który chce zrozumieć siebie do końca - nie wedle jakiś tylko doraźnych, częściowych, czasem powierzchownych,  a nawet pozornych kryteriów i miar swojej własnej istoty - musi ze swoim niepokojem, niepewnością, a także słabością i grzesznością, ze swoim życiem i śmiercią, przybliżyć się do Chrystusa. Musi niejako w i Niego wejść z sobą samym,  musi sobie przyswoić, zasymilować całą rzeczywistość Wcielenia i Odkupienia,  aby siebie odnaleźć.
\end{quote}
\begin{itemize}
\item {} 
„Wejść całym sobą”, „zasymilować całą rzeczywistość Wcielenia i Odkupienia” - to niesamowicie wymagające słowa. Czy odczuwamy strach, niepokój słysząc tak radykalne zdania? Dlaczego?

\item {} 
Czy mamy doświadczenie takiego otwartego  zbliżenia się do Chrystusa całym sobą?

\end{itemize}

Jest to to o czym mówiliśmy wczoraj: człowiek który odkrył już swój niepokój, niepewność, słabość, grzeszność nadal nie potrafi zrozumieć własnego sensu życia, ma więcej pytań niż odpowiedzi. Na tym właśnie etapie trzeba szczególnie zwrócić się do Jezusa... Tylko on da nam możliwość poszukiwania odpowiedzi na nasze pytania głębiej.  Nie  należy się tego wstydzić,  że robimy to po to aby „poznać siebie”.


\section{Zaduma  nad samym sobą}
\label{babice2006-wiosna-knurow/spotkanie2:zaduma-nad-samym-soba}\begin{quote}

Jeśli  dokona się w i człowieku ów dogłębny  proces, wówczas owocuje on nie tylko uwielbieniem Boga, ale także głębokim zdumieniem nad sobą samym. Jakąż wartość musi mieć w i oczach Stwórcy człowiek, skoro zasłużył na takiego i tak potężnego Odkupiciela, skoro Bóg Syna swego Jednorodzonego   dał, ażeby on, człowiek, nie zginął, ale miał życie wieczne.
\end{quote}
\begin{itemize}
\item {} 
W jaki sposób my sami troszczymy się o to, aby ten proces inicjować?

\item {} 
Czy potrafimy dostrzec w i sobie tą wielką wartość? Czy nie jest tak, że często krytykując sami siebie odbieramy sobie tą godność? (tzn. robimy z siebie ludzi „za małych“)

\end{itemize}


\section{Odważne przyjęcie Miłości}
\label{babice2006-wiosna-knurow/spotkanie2:odwazne-przyjecie-milosci}
Boża miłość, jej ogrom, może nas czasami przytłaczać. Tak jak mówiliśmy wczoraj w i czasie konferencji wieczornej o listach. Intencją tego który obdarza miłością nie jest to abyśmy czuli się przytłoczeni. Ciężko jest przyjmować miłość.
\begin{itemize}
\item {} 
Czy doświadczamy tego cierpienia związanego z \textbf{przyjmowaniem miłości?} Co robimy, aby cierpieniu temu nadawać sens?

\end{itemize}

Cierpienie  jest  w i miłość  wpisane (przynajmniej  tą  tutaj na ziemi). Uciekanie przed tym cierpieniem implikuje uciekanie przed przyjęciem pełni miłości. W i życiu każdego człowieka musi nastąpić moment trudnej decyzji, aby zacząć kochać (otworzyć się na to, aby ktoś nas kochał).
\begin{itemize}
\item {} 
Jak my radzimy sobie z tą decyzją? Czy potrafimy przyjąć taką naszą sytuację z optymizmem, pogodą ducha?

\end{itemize}

Przeczytajmy:
\begin{quote}

Każda prawdziwa Miłość musi mieć swój Wielki Piątek.

\begin{flushright}
---kard. Stefan Wyszyński
\end{flushright}
\end{quote}
\begin{itemize}
\item {} 
Jak rozumiemy te słowa kardynała?

\end{itemize}

Nie przez przypadek miłowanie porównuje się do palącej  się świecy - spalamy siebie dla drugich. Dochodzimy do najważniejszego pytania:
\begin{itemize}
\item {} 
\textbf{Czy potrafimy  bardziej radować się światłem które dzięki temu otrzymują drudzy niż smucić z powodu naszych niedogodności?}

\end{itemize}


\section{Zastosowanie}
\label{babice2006-wiosna-knurow/spotkanie2:zastosowanie}
Metoda Ruchu Światło-Życie nie pozwala na to, abyśmy pozostawali tylko na poziomie wiedzy i i informacji. Każde spotkanie powinno zakończyć się
„zastosowaniem”, czyli konkretnym wprowadzeniem w i życie treści.

Propozycja, aby każdy z i nas podjął teraz jedno małe postanowienie w i związku z i omawianym tematem, które liczy, że będzie mógł zrealizować jeszcze na tych rekolekcjach.


\section{Modlitwa końcowa}
\label{babice2006-wiosna-knurow/spotkanie2:modlitwa-koncowa}
Dowolna w i intencji dobrego przeżycia Triduum Paschalnego.


\part{Eucharystia}
\label{babice2006-jesien-knurow/index::doc}\label{babice2006-jesien-knurow/index:eucharystia}\paragraph{Podstawowe informacje o rekolekcjach}
\begin{itemize}
\item {} \begin{description}
\item[{Data:}] \leavevmode\begin{itemize}
\item {} 
Jesień 2006 roku

\end{itemize}

\end{description}

\item {} \begin{description}
\item[{Miejsce pierwszych rekolekcji:}] \leavevmode\begin{itemize}
\item {} 
Babice, diecezja gliwicka

\end{itemize}

\end{description}

\item {} \begin{description}
\item[{Miejsce powstania konspektów:}] \leavevmode\begin{itemize}
\item {} 
Knurów, archidiecezja katowicka

\end{itemize}

\end{description}

\item {} \begin{description}
\item[{Grupa docelowa:}] \leavevmode\begin{itemize}
\item {} 
Oaza parafialna, od preoazy do formacji przed III st.

\end{itemize}

\end{description}

\item {} \begin{description}
\item[{Opis rekolekcji:}] \leavevmode\begin{itemize}
\item {} 
BRAK

\end{itemize}

\end{description}

\end{itemize}


\chapter{Spotkanie 1 - Eucharystia jako spotkanie z człowiekiem i Bogiem}
\label{babice2006-jesien-knurow/spotkanie1::doc}\label{babice2006-jesien-knurow/spotkanie1:spotkanie-1-eucharystia-jako-spotkanie-z-czlowiekiem-i-bogiem}

\section{Cel Spotkania}
\label{babice2006-jesien-knurow/spotkanie1:cel-spotkania}
Pokierować ludzi do spojrzenia na Eucharystie jako na Spotkanie. Uświadomienie uczestnikom wartości spotkania samego w sobie. Refleksja nad naszym przygotowaniem do spotkania.


\section{Wprowadzenie dla Animatora}
\label{babice2006-jesien-knurow/spotkanie1:wprowadzenie-dla-animatora}
Takie słowa do przemyślenia. śeby zdać sobie sprawę czym jest spotkanie. W starszych grupach być może warto przeczytać je jako wprowadzenie.
\begin{quote}

Dobro ma na świecie swoje imię, a nie nazwę. Jest: Adamem, Ewą, Abrahamem (...)” — pisze ksiądz Tischner. Lecz co to dokładnie znaczy? Jeśli Dobro jest czymś tak radykalnie nie-przedmiotowym, czymś, co nie ma nazwy, czymś, co występuje tylko pod postacią konkretnych ludzi, to jak możemy je mierzyć? Jak właściwie można oceniać, czy dany czyn jest dobry, lepszy, gorszy? Gdzie jest miara, do której powinniśmy odnieść postępowanie Adama, Abrahama, Ewy? I dalej: co łączy Adama, Abrahama, Ewę? Widząc człowieka, zdaje się nam sugerować ksiądz Tischner, powinniśmy doznawać olśnienia, że oto przed nami pojawił się ktoś nieskończenie godny tego, by go miłować.
\end{quote}

Spotkanie to nie będzie akcentowało się na słowie „Eucharystia” mimo wszystko. Będzie na jej temat jeszcze wiele w sobotę oraz niedziele. Chcemy w tym spotkaniu się przygotować. Zastanowic się nad istotą samego spotkania (z Bogiem lub z Człowiekiem – generalnie OSOBĄ…..). Syntezą tego spotkania będzie: Eucharystia to miejsce podwójnego spotkania.


\section{Modlitwa}
\label{babice2006-jesien-knurow/spotkanie1:modlitwa}
Dostosowana do możliwości grupy.


\section{Refleksja nad „Spotkaniem”}
\label{babice2006-jesien-knurow/spotkanie1:refleksja-nad-spotkaniem}
Najpierw przeczytajmy takie zdanie. Być może będziemy się z nim zgadzać lub nie – będzie ono wprowadzeniem do naszej dalszej rozmowy.
\begin{quote}

Niewątpliwie zdecydowana większość tego, co dziś między ludźmi nazywa się rozmową, należałoby trafniej, w dokładnym sensie tego słowa, określić jako gadanina. Na ogół ludzie nie mówią naprawdę do siebie; każdy, chociaż zwrócony do drugiego, mówi w rzeczywistości do fikcyjnej instancji, której istnienie wyczerpuje się w wysłuchaniu go.

\begin{flushright}
---Martin Buber
\end{flushright}
\end{quote}

Warto zadać sobie pytanie nad jakością naszych spotkań (wszelkich). Spróbować odnaleźć w swoim życiu takie prawdziwe rozmowy/spotkania. Warto sobie je wypisać na jakiejś kartce (każdy dla siebie).
\begin{itemize}
\item {} 
Czy mamy doświadczeń takiego prawdziwego spotkania się z drugą osobą?

\item {} 
Czy możemy się podzielić w jakich to było okolicznościach?

\end{itemize}

Spotkanie to dar. Coś absolutnie niesamowitego. Jeżeli mamy niewiele wpisów na naszych kartkach to doskonale to rozumiemy bo prawdopodobnie „tęsknimy” za prawdziwą rozmową. Jeżeli nasze kartki są zapełnione to prawdopodobnie doświadczyliśmy tego i cenimy sobie te momenty.
\begin{quote}

Dąży się dzisiaj nie tylko do „odczarowania” kontaktów międzyludzkich – z czym można by się nawet zgodzić – lecz także do ich radykalnego odtajemniczenia. To, co osobowe, nieustannie bliskie misterium, niegdyś podstawa najcichszych uniesień, ulega niwelacji.

\begin{flushright}
---Martin Buber
\end{flushright}
\end{quote}
\begin{itemize}
\item {} 
Czy zgadzamy się z tym zdaniem? Czy ``tęsknimy'' za taka prawdziwa rozmowa/spotkaniem? Czy potrafi nam tego brakować?

\item {} 
Czy podziękowaliśmy komuś za spotkanie? Czy nie jest tak, że wobec wielu osób nie wyraziliśmy swojej wdzięczności za rozmowę?

\end{itemize}

Postanowienie: Świetną okazją, aby „podziękować” za ten „cud spotkania” osobą które pojawiły się w naszym życiu jest jutrzejsza Eucharystia! Niech to będzie pierwszym owocem tego spotkania (o ile ktoś chce) – aby jutrzejszą Eucharystię ofiarować w ich intencji.


\section{Cechy Prawdziwego Spotkania}
\label{babice2006-jesien-knurow/spotkanie1:cechy-prawdziwego-spotkania}
Powstaje pytanie:
\begin{itemize}
\item {} 
Dlaczego tak niewiele Rozmowy/Spotkań jest w naszym życiu?

\item {} 
Dlaczego jest to niewygodne dla świata?

\end{itemize}

Chciałoby się zapytać: co przeszkadza nam w Spotkaniu. Zastanowimy się jakie warunki są potrzebne, aby mogło dojść do Spotkania.

Animator będzie pokazywał poszczególne „wymogi” (odczytuje, a następnie kładzie kartkę na środek! Punkty do „pocięcia nożyczkami” znajdują się na ostatniej stronie konspektu) a uczestnicy będą przekuwać je na język codzienności i konkretu życia – mogą podawać przykłady, wytłumaczyć własnym językiem (powinni sami sobie tłumaczyć). śeby lepiej zrozumieć każdy punkt niech przedstawią sytuację w której w Spotkaniu zapominamy o danej rzeczy (jakie zdania wtedy mogą zostać wypowiedziane, jakie sytuacje zaistnieć)

Uczestnicy mogą rozszerzyć tą listę!!!
\begin{itemize}
\item {} 
Cechy Prawdziwego Spotkania:
\begin{itemize}
\item {} 
Zwrócenie się do partnera w całej prawdzie. Mówię drugiemu jako osobie: „tak” - potwierdzamy jego osobę

\item {} 
Każdy musi włączyć samego siebie. Każdy zawsze mówi to, co myśli jednak: „otwartość jest dokładnym przeciwieństwem gadaniny.”

\item {} 
Nie kieruje się myślą o własnym oddziaływaniu jako mówcy tego, co powinno przez niego przemówić! Gdy zamiast tego, co winno być powiedziane, chce zaprezentować swoje samolubne „ja” staje się niszczycielem

\item {} 
Rozmowa ze swojej natury jest płodna. Rozmowa owocuje…tworzy w nas kogoś innego, odkrywamy sami siebie w niej bardziej...

\item {} 
Głównym warunkiem zaistnienia prawdziwej rozmowy jest to, by każdy uważał swego partnera za tą, właśnie tą osobę. Postrzegam ją, postrzegam, że jest ona inna, istotowo inna niż ja, w ten określony, jej tylko właściwy i niepowtarzalny sposób istotowo inna niż ja, i akceptuję tą osobę, którą postrzegłem, tak że mogę z całą powagą skierować do niej, właśnie do niej, moje słowa.

\item {} 
W kontaktach międzyludzkich dominuje dziś spojrzenie analityczne, redukujące i wywodzące. Jest analityczne – lub raczej pseudoanalityczne – ponieważ traktuje cały byt cielesno-duchowy jako złożony, a zatem i rozkładalny;(...) Spojrzenie jest redukcyjne, ponieważ różnorodność osoby (...) chce sprowadzić do schematycznych i powtarzających się wszędzie struktur.

\end{itemize}

\end{itemize}

\textbf{Te DWA ostatnie cytaty animator KONIECZNIE musi zmienić jeżeli uzna to za stosowne do grupy!!!}


\section{W poszukiwaniu Spotkania}
\label{babice2006-jesien-knurow/spotkanie1:w-poszukiwaniu-spotkania}
Wiemy, że coraz rzadziej naprawdę się ze sobą spotykamy. Wiemy czym się takie spotkania powinny charakteryzować i co musi być spełnione, aby mogło do nich dojść. Czy to jednak nam wystarcza, aby móc ich doświadczyć? Na pewno jest to pomocne, może nawet niezbędne – więc warto było. Jednak i tak wiele z nas chodzi po świecie z tęsknotą za Prawdziwym Spotkaniem.

Człowiek jest tak stworzony. Poszukujemy… Im dłużej na świecie żyjemy tym bardziej człowiek doświadcza tego, że tego jego różnego rodzaju spotkania nie dotykają pełni (lub nawet nie dotykają powierzchni?). Zbliżamy się w tym momencie do tematu naszych rekolekcji…do Eucharystii….odpowiedzi Boga na te nasze ogromne pragnienie Prawdziwego Spotkania….
\begin{itemize}
\item {} 
Czy spojrzeliście kiedykolwiek w ten sposób na Eucharystię?

\item {} 
Czy kiedy odczuwaliście tą swoistego rodzaju „samotność duszy” pomyśleliście, że Eucharystia jest właśnie odpowiedzią? (nie wiem czy nie jest zbyt osobiste)

\end{itemize}

\emph{Tak naprawdę w tym momencie może nawet powinniśmy zatrzymać spotkanie. Nawet nie odpowiadać na te pytania głośno. Zostać w takim zdumieniu – jest coś co odczuwam
i czego szukam…Bóg wie o tym co głęboko we mnie siedzi….wiedział już zanim mnie stworzył i tak przygotował Plan Zbawienia, że dał nam Eucharystie...Po prostu Miłość.}


\section{Umocnienie w odkrywaniu i zastosowanie}
\label{babice2006-jesien-knurow/spotkanie1:umocnienie-w-odkrywaniu-i-zastosowanie}
\emph{(to nie jest cześć spotkania „na koniec” tylko bardzo ważny punkt któremu trzeba przewidzieć odpowiednią ilość czasu jeżeli spotkanie ma mieć sens większy niż ćwiczenia z tekstem!)}

Wiadomo jak jest - Eucharystia jest trudna. śycie toczy się i nie jest łatwo myśleć w takich kategoriach o Mszy Świętej szczególnie jeżeli coś innego w tym momencie jest do zrobienia. Czasami mamy po prostu wiele ludzkich wątpliwości co do tego co dzieje się na Ołtarzu. To normalne. Te rekolekcje są okazją do próby podjęcia jakiegoś wysiłku, aby spróbować pogłębić własne przeżywanie i patrzenie na Eucharystie. Chcemy, aby Msza Święta była dla nas właśnie takim upragnionym Spotkaniem. Jak będziemy się starać o to?

Tworzymy grupę…..grupę która się właśnie SPOTYKA ze sobą. Możemy sobie pomóc. Nie będziemy teraz wyliczać różnych porad „z podręczników”! Stwórzmy sami taką listę rzeczy która będzie miała nam pomóc.
\begin{enumerate}
\item {} 
\textbf{Każdy uczestnik wypisuje na jednej wspólnej kartce rzeczy (postanowienia) do których się zobowiązuje, aby pomóc wzajemnie sobie przeżyć dobrze ten czas pod kątem odkrycia Eucharystii. (spotkanie z Człowiekiem)}

\item {} 
\textbf{Każdy uczestnik wypisuje na własnej kartce rzeczy (postanowienia) do których się zobowiązuje osobiście, aby dobrze ten czas pod kątem odkrycia Eucharystii. (spotkanie z Bogiem)}

\end{enumerate}

Kartki „dla siebie” chowamy w takie miejsce aby było zawsze „pod ręką”….nie zakopujemy tam gdzie nigdy już do niego nie zajrzymy!

Kartkę „wobec grupy” zachowujemy i rozkładamy przed każdym spotkaniem koło świeczki – tak abyśmy sobie przypominali co najmniej raz na dzień nasze postanowienia.

Pamiętajmy o intencji jutrzejszej Mszy!


\chapter{Spotkanie 2 - Eucharystia jako pokarm}
\label{babice2006-jesien-knurow/spotkanie2:spotkanie-2-eucharystia-jako-pokarm}\label{babice2006-jesien-knurow/spotkanie2::doc}

\section{Wprowadzenie dla Animatora}
\label{babice2006-jesien-knurow/spotkanie2:wprowadzenie-dla-animatora}
Celem dzisiejszego spotkania jest:
\begin{enumerate}
\item {} 
zwrócenie uwagi na chleb, który jest podstawowym pokarmem i wyjaśnienie dlaczego właśnie pod tym symbolem Jezus zostawił nam siebie;

\item {} 
ukazanie sytuacji w Starym Testamencie, znaczenie chleba, pokarmu od Boga dla Izraelitów;

\item {} 
uświadomienie jakie zmiany do tej relacji wprowadza Jezus Chrystus, który jest \textbf{prawdziwym} pokarmem;

\item {} 
zadanie pytania o rolę Eucharystii, przyjęcia Jezusa w komunii, w moim życiu i tym samy pogłębienie świadomości tego momentu Eucharystii.

\end{enumerate}

Dzisiejsze spotkanie odbywa się późno, jesteśmy już po Namiocie Spotkania oraz pieczeniu chleba. Do tych dwóch punktów to spotkanie będzie się przede wszystkim odwoływało. Jest to spotkanie w dużej mierze oparte o fragmenty z Pisma Św. Nie chodzi tu jednak, aby szybko je przeczytać i przejść jak najszybciej do następnego fragmentu. Każdy fragment należy bardzo dokładnie z uczestnikami omówić, wyjaśnić. Pomocne, zwłaszcza w grupach młodszych będzie opowiedzenie sytuacji ukazanej w Piśmie Św., tj. gdzie się odbywa, kto bierze udział, jak to mogło wyglądać, jakie towarzyszyły temu emocje itd., odwołanie się do znaczenia chleba w ich życiu, nawet poprzez takie pytania, jak: „jak często jecie chleb, jaki lubicie najbardziej itp.” oraz odwołanie do pieczenia chleba, które miało miejsce wcześniej. „Jak to jest być piekarzem?” Generalnie należy dążyć, aby to spotkanie było jak najbardziej ROZMOWĄ z uczestnikami. Absolutnie nie może być to wykład teologiczny aniamtora. W czasie spotkania należy ukazać różnice pomiędzy ST a NT, co nowego wnosi Jezus Chrystus.


\section{Wstęp}
\label{babice2006-jesien-knurow/spotkanie2:wstep}
Przeczytajmy:
\begin{quote}

Następnie wyruszyli z Elim. Przybyło zaś całe zgromadzenie Izraelitów na pustynię Sin, położoną między Elim a Synajem, piętnastego dnia drugiego miesiąca od ich wyjścia z ziemi egipskiej. I zaczęło szemrać na pustyni całe zgromadzenie Izraelitów przeciw Mojżeszowi i przeciw Aaronowi. Izraelici mówili im: «Obyśmy pomarli z ręki Pana w ziemi egipskiej, gdzieśmy zasiadali przed garnkami mięsa i jadali chleb do sytości! Wyprowadziliście nas na tę pustynię, aby głodem umorzyć całą tę rzeszę». Pan powiedział wówczas do Mojżesza: «Oto ześlę wam chleb z nieba, na kształt deszczu. I będzie wychodził lud, i każdego dnia będzie zbierał według potrzeby dziennej. Chcę ich także doświadczyć, czy pójdą za moimi rozkazami czy też nie. Lecz w dniu szóstym zrobią zapas tego, co przyniosą, a będzie to podwójna ilość tego, co będą zbierać codziennie». Mojżesz i Aaron powiedzieli do społeczności Izraelitów: «Tego wieczora ujrzycie, że to Pan wyprowadził was z ziemi egipskiej. A rano ujrzycie chwałę Pana, gdyż usłyszał On, że szemrzecie przeciw Panu. Czymże my jesteśmy, że szemrzecie przeciw nam?» Mojżesz powiedział: «Wieczorem Pan da wam mięso do jedzenia, a rano chleb do sytości, bo słyszał Pan szemranie wasze przeciw Niemu. Czymże bowiem my jesteśmy? Nie szemraliście przeciwko nam, ale przeciw Panu!» Mojżesz rzekł do Aarona: «Powiedz całemu zgromadzeniu Izraelitów: Zbliżcie się do Pana, gdyż słyszał wasze szemrania». W czasie przemowy Aarona do całego zgromadzenia Izraelitów spojrzeli ku pustyni i ukazała się im w obłoku chwała Pana. I przemówił Pan do Mojżesza tymi słowami: «Słyszałem szemranie Izraelitów. Powiedz im tak: O zmierzchu będziecie jeść mięso, a rano nasycicie się chlebem. Poznacie wtedy, że Ja, Pan, jestem waszym Bogiem». Rzeczywiście wieczorem przyleciały przepiórki i pokryły obóz, a nazajutrz rano warstwa rosy leżała dokoła obozu. Gdy się warstwa rosy uniosła ku górze, wówczas na pustyni leżało coś drobnego, ziarnistego, niby szron na ziemi. Na widok tego Izraelici pytali się wzajemnie: «Co to jest?» - gdyż nie wiedzieli, co to było. Wtedy powiedział do nich Mojżesz: «To jest chleb, który daje wam Pan na pokarm.  To zaś nakazał wam Pan: Każdy z was zbierze dla siebie według swej potrzeby, omer na głowę. Każdy z was przyniesie według liczby osób, które należą do jego namiotu». Izraelici uczynili tak i zebrali jedni dużo, drudzy mało. Gdy mierzyli swój zbiór omerem, to ten, który zebrał wiele, nie miał nic zbywającego, kto zaś za mało zebrał, nie miał żadnego braku - każdy zebrał według swych potrzeb. {[}...{]} Izraelici jedli mannę przez czterdzieści lat, aż przybyli do ziemi zamieszkałej. Jedli mannę, aż przybyli do granic ziemi Kanaan.

\begin{flushright}
---Wj 16, 1-18.35
\end{flushright}
\end{quote}

Izraelici szemrają przeciw Panu, gdyż nie mają pokarmu, nie mają chleba. Uważają, że lepiej dla nich byłoby żyć w Egipcie, gdzie przynajmniej mieli zagwarantowany pokarm. Przy tej okazji
warto porozmawiać z uczestnikami o tym, jak wygląda chleb, z czego jest robiony. W Izrealu chleb (podpłomyk) był robiony z wody, soli i mąki (ew. jeszcze oliwy). Są to podstawowe składniki. Najprostsze. Chleb przaśny, używany podczas Paschy był robiony bez zakwasu (czyli reszty ciasta pozostałej z poprzedniego wypieku), tylko z wody i mąki. Po pieczeniu chleba uczestnicy powinni mieć świadomość prostoty chleba. Niech na początku tego spotkania podzielą się doświadczeniami, przemyśleniami, jakich doznali podczas pieczenia chleba. Poprzez tą prostotę chleb jest podstawowym pożywieniem wielu ludzi, zwłaszacza ubogich (gdzie jednocześnie jest jedynym składnikiem pożywienia). Zapytajmy o co najczęściej proszą ludzie ubodzy (oprócz pieniędzy).

Następnie porozmawiajmy o tym, jak traktowany jest chleb w naszych domach.
\begin{itemize}
\item {} 
Jak dużo chleba w naszych domach jest marnowane? Jakie w naszych domach panuje podejście do chleba? Czy podchodzimy z szacunkiem do chleba (tradycja b. rozpowszechniona na Śląsku) ?

\end{itemize}


\section{Bóg troszczy się o nas}
\label{babice2006-jesien-knurow/spotkanie2:bog-troszczy-sie-o-nas}
Pan Bóg troszczy się także o nasze życie codzienne, o nasz pokarm doczesny (w „Ojcze nasz” prosimy o chleb powszedni, który może być rozumiany także jako ten ziemski). Bóg wysłuchuje szemrania Izraelitów, zsyła swój pokarm. Przy tym fragmencie z księgi Wyjścia zwróćmy uwagę na dwie rzeczy. Po pierwsze Izraelici pytali „co to jest?”, było to dla nich zjawisko niezwykłe, nowe. Po drugie każdy mógł zebrać tyle, żeby najeść się do syta. Bóg nie pozwala nikomu być głodnym. Porozmawiajmy o tym fragmencie, co się w nim wydarzyło, jeki były przyczyny, jakie skuki. Niech to będzie luźna rozmowa na temat tego fragmentu. Można przytoczyć fragment psalmu:
\begin{quote}

\begin{DUlineblock}{0em}
\item[] Wystawiali Boga na próbę w swych sercach, żądając strawy dla swego pożądania.
\item[] Mówili przeciw Bogu, rzekli: «Czyż Bóg potrafi nakryć stół w pustyni?» Oto w skałę uderzył,
\item[] popłynęły wody i wytrysnęły strumienie:
\item[] «Czy także potrafi dać chleba
\item[] albo ludowi swemu przygotować mięso?» Toteż gdy Pan usłyszał, zapalił się gniewem, i ogień rozgorzał przeciw Jakubowi,
\item[] \textless{}i gniew jeszcze powstał przeciw Izraelowi\textgreater{}, że nie uwierzyli w Boga
\item[] i nie zaufali Jego pomocy.
\item[] Potem z góry wydał rozkaz chmurom i bramy nieba otworzył,
\item[] i spuścił jak deszcz mannę do jedzenia:
\item[] dał im zboże z nieba.
\item[] Człowiek chleb mocarzy spożywał - żywności zesłał im do syta.
\end{DUlineblock}

\begin{flushright}
---Ps 78, 18-25
\end{flushright}
\end{quote}

Zapytajmy:
\begin{itemize}
\item {} 
Dlaczego Bóg zsyła Izrealitom pokarm?

\item {} 
Jaka jest reakcja Izrealitów na to wydarzenie?

\item {} 
Jak zachowują się wobec Boga po tym wydarzeniu? Czy ich postępowanie zmienia się?

\item {} 
Co znaczy określenie „chleb mocarzy”? Jak je rozumiesz?

\item {} 
Powracając do fragmentu z księgi Wyjścia: jak długo Izraelici spożywali chleb od Pana? Dlaczego? (Odkryć znaczenie symboliczne: wędrówka – całe nasze życie, Kanaan – niebo, nie wspominajmy jeszcze o Ciele Jezusa Chrystusa.)

\end{itemize}

Bóg karmi swój naród przez całą jego wedrówkę, przez całą drogę. Pokarm ten pozwala im iść do celu, jaki przeznaczył im Bóg. Jednak ten cud nie powoduje, że Izraelici przestają szemrać, wątpić. Dlatego większość z nich pomarła i nie doszła do ziemi obiecanej.


\section{Sakrament}
\label{babice2006-jesien-knurow/spotkanie2:sakrament}
Przeczytajmy:
\begin{quote}

\begin{DUlineblock}{0em}
\item[] „Na świecie było {[}Słowo{]}, a świat stał się przez Nie, lecz świat Go nie poznał.
\item[] Przyszło do swojej własności, a swoi Go nie przyjęli.
\item[] Wszystkim tym jednak, którzy Je przyjęli, dało moc, aby się stali dziećmi Bożymi, tym, którzy wierzą w imię Jego -
\item[] którzy ani z krwi,
\item[] ani z żądzy ciała, ani z woli męża,
\item[] ale z Boga się narodzili. A Słowo stało się ciałem
\item[] i zamieszkało wśród nas.
\item[] I oglądaliśmy Jego chwałę,
\end{DUlineblock}

'' chwałę, jaką Jednorodzony otrzymuje od Ojca, pełen łaski i prawdy.

\begin{flushright}
---J 1, 10-14
\end{flushright}
\end{quote}

W oparciu o powyższy fragment oraz Namiot Spotkania porozmawiajmyo Eucharystii, o Chlebie jaki możemy dziś spożywać. Niech uczestnicy podzielą się swoimi przemyśleniami z Namiotu Spotkania.
\begin{itemize}
\item {} 
Jakie są analogie pomiędzy naszym życiem, a wędrówką Izraelitów? Zwróćmy uwagę, że manna była odpowiedzią Boga na szemranie ludzi. Scana z Namiotu Spotkania odbywa się po rozmnożeniu chleba, kiedy Jezus wyrzuca ludziom, że idą za nim tylko dlatego, że najedli się do syta, natomiast nie troszczą się o życie wieczne. Jaka następuje tutaj różnica pomiędzy Starym Testemantem a Nowym? Zastanówmy się jak jest w naszym życiu – czy przystępujemy do komunii tylko dlatego, żeby się najeść do syta?

\end{itemize}

Nawiążmy do przemyśleń z Namiotu Spotkania:
\begin{itemize}
\item {} 
Czym tak naprawdę jest przyjęcie Jezusa, spożywanie Jego Ciała i Krwi? Jest to o wiele większy dar niż tylko pomoc w doczesnym życiu jak to miało miejsce w ST. Tutaj chleb oznacza życie wieczne, oznacza zjednoczenie z Chrystusem (o czym jutro, więc nie rozwijajmy tego wątku), oznacza wreszcie UDZIAŁ W OFIERZE CHRYSTUSA, a także zadatek, przedsmak nieba.

\end{itemize}


\section{„Co to jest?” = „lecz świat Go nie poznał”}
\label{babice2006-jesien-knurow/spotkanie2:co-to-jest-lecz-swiat-go-nie-poznal}
Bóg pragnął pozostać między nami, nie chciał zostawić nas sierotami - „A Słowo stało się ciałem i zamieszkało wśród nas.” Zastanówmy się przez chwilę dlaczego Bóg pozostał z nami, dlaczego pozwolił nam spożywać się.
\begin{quote}

Pan nasz Jezus Chrystus chciał zostawić nam zbawienie w swoim Ciele i Krwi. Dlaczego zaś każe nam spożywać swoje Ciało i Krew? Bo jest pokorny. Gdyby nie był pokorny, nie moglibyśmy Go spożywać ani pić. Patrz na Jego nieskończoność: „Na początku było Słowo i Słowo było u Boga i Bogiem było Słowo” (J 1,1). Oto Pokarm wiekuisty: lecz spożywają Go aniołowie, spożywają moce najwyższe i duchy niebieskie. Spożywając, nabierają mocy, a Ten, który obdarza ich nasyceniem i radością, pozostaje nie naruszony. Otóż, co człowiekowi do tego Pokarmu? Kto znajdzie serce dostosowane do takiego Pokarmu? Zatem w mleko się On przemienił i dostosował do maluczkich. Oto w jaki sposób Mądrość Boża karmi nas tym Pokarmem: „Słowo ciałem się stało i zamieszkało wśród nas” (J 1,14). Patrzcie, co za pokora! „Dałeś im Chleb z nieba, Chleb aniołów spożywa człowiek” (Ps 77,24). Słowo Przedwieczne, którym żywią się aniołowie, równe Ojcu, jest spożywane przez człowieka. Ponieważ „istniejąc w postaci Bożej, nie skorzystał ze sposobności, aby być na równi z Bogiem” — takim właśnie żywią się aniołowie — „lecz ogołocił samego siebie” — aby człowiek mógł spożywać Chleb aniołów — „przyjąwszy postać sługi, stawszy się podobnym do ludzi. A w zewnętrznym przejawie uznany za człowieka, uniżył samego siebie, stawszy się posłusznym aż do śmierci, i to śmierci na krzyżu” (Flp 2,6—8) — abyśmy z krzyża otrzymali Ciało i Krew Pańską jako nową ofiarę.

\begin{flushright}
---św. Augustyn
\end{flushright}
\end{quote}

Jezus pragnie pozostać z nami, gdyż jest to przejaw Jego miłości do nas. On tak nas umiłował, że pragnie, abyśmy mogli być z nim jedno, abyśmy mogli już tu na ziemi dostąpić, chociaż częściowo, tego, co nas czeka w niebie. Chrześcijanin to człowiek wypatrujący powtórnego przyjścia swojego Zbawiciela, człowiek wpatrzony w niebo. Aby ten czas oczekiwania był bardziej znośny dla człowieka, Jezus zostawił nam swoje Ciało, abyśmy potrafili czekać.
\begin{itemize}
\item {} 
Czy współczesny człowiek pragnie zjednoczenia z Bogiem? Czy jest to także nasze największe pragnienie?

\end{itemize}

Nie może być Eucharystia źródłem i szczytem życia chrześciajanina, jego najważniejszym momentem, jeśli nie pragnie on przede wszystkim zjednoczenia ze swoim Zbawcą, jeśli nie będzie to jego najważniejszym celem. Stawiamy sobie tutaj pytanie o to, czy Bóg jest dla mnie najważniejszy. Jeśli popatrzymy na tą sytuację od strony Boga dostrzeżemy, że dla Niego nie ma nic ważniejszego jak zjednoczyć się ze swoim ludem, ze swoim stworzeniem. On ogołocił samego siebie, aby móc być z nami!
\begin{quote}

W Eucharystii, w odróżnieniu od każdego innego sakramentu, tajamnica {[}komunii{]} jast tak doskonała, iż prowadzi do szczytu wszelkich dóbr: to tu znajduje kres wszelkie pragnienie, ponieważ tu otrzymujemy Boga i Bóg wchodzi w doskonałe zjednoczenie z nami.

\begin{flushright}
---Jan Paweł II
\end{flushright}
\end{quote}

\textbf{To właśnie znaczy, że Eucharystia jest szczytem życia chrześcijanina, ponieważ w niej dokonuje się największe pragnienie chrześcijanina: jednoczy się z Bogiem.}


\section{Podsumowanie}
\label{babice2006-jesien-knurow/spotkanie2:podsumowanie}\begin{quote}

Posłuchajcie o małej dziewczynce, którą znalazłam na jednej z ulic Kalkuty. Miała może sześć lat, a po jej twarze poznałam, że jest głodna i nie jadła od kilku dni. Dałam jej kromkę chleba, a ona zaczęła ją jeść, powoli, po jednym okruszku. Powiedziałam: Jedz chleb, nie bój się, jedz. Dziecko powiedziało: Boję się, bo jak chleb się skończy, znowu będę głodna

\begin{flushright}
---Matka Teresa z Kalkuty
\end{flushright}
\end{quote}

Człowiek nie może żyć bez codziennego pokarmu, chrześcijanin nie może żyć bez Ciała i Krwi Chrystusa. Spożywając je już nie my żyjemy, ale żyje w nas Chrystus.
\begin{itemize}
\item {} 
Czy w nas jest także obawa, że jak chleb się skończy, to znów będziemy głodni?

\item {} 
Czy mieliśmy taki okres w naszym życiu, kiedy brakowało nam tego Pokarmu? Jak czujemy się, kiedy nie możemy przyjąć Jezusa do naszego serca?

\item {} 
Czy mogę zrobić coś, żeby sprawić, aby tych okresów było jak najmniej?

\end{itemize}


\section{Zastosowanie}
\label{babice2006-jesien-knurow/spotkanie2:zastosowanie}\begin{itemize}
\item {} 
postaram się zmienić moje podejście do pokarmu, jedzenia (o ile to konieczne);

\item {} 
postaram się uniknąć w ciągu najbliższych tygodni/miesiąca okresów, w których nie będę mógł/mogła przystąpić do komunii św. (rozwiązaniem tej sytuacji jest jak najszybsze przystąpienie do sakramentu pokuty, a nie np. unikanie mszy...)

\end{itemize}


\chapter{Spotkanie 3 - Eucharystia jako zjednoczenie z ludźmi i Bogiem}
\label{babice2006-jesien-knurow/spotkanie3:spotkanie-3-eucharystia-jako-zjednoczenie-z-ludzmi-i-bogiem}\label{babice2006-jesien-knurow/spotkanie3::doc}

\section{Wprowadzenie dla Animatora}
\label{babice2006-jesien-knurow/spotkanie3:wprowadzenie-dla-animatora}
Celem dzisiejszego spotkania jest:
\begin{enumerate}
\item {} 
zwrócenie uwagi na potrzebę przygotowania się do Eucharystii;

\item {} 
ukazanie jedności ludzi przez pryzmat przyjęcia Ciała i Krwi Chrystusa;

\item {} 
ukazanie Eucharystii jako źródło postaw, które powinniśmy praktykować w naszym codziennym życiu.

\end{enumerate}

Dzisiejsze spotkanie odbywa się po niedzielnej Eucharystii, dlatego wnioski i postanowienia dotyczyć będą każdej następnej mszy św. po powrocie do domu, bez wspólnoty rekolekcyjnej. Dlatego najważniejszą rzeczą jest, abyśmy uświadomili uczestników na te nowe wspólnoty, często przypadkowych ludzi, abyśmy przede wszystkim pokazali na dzisiejszym spotkaniu jedność i równość ludzi uczestniczących we mszy św. i karmiących się Ciałem i Krwią Chrystusa. Drugim ważnym punktem jest ukazanie mszy św. jako źródła postaw i wzór do naśladowania w naszym codziennym życiu, ukazanie Eucharystii, gdzie ma miejsce doskonała wspólnota ludzi w Bogu, którą powinniśmy starać się tworzyć również poza kościołem.

Dziś najważniejsze jest odwołanie do przeżyć, odczuć uczesntników. Niech opowiedzą jak przeżyli te msze podczas tych rekolekcji, czy różniły się one od innych dotychczasowych przez nich przeżywanych. Ta reflekcja może przebiegać wg pytań w konspekcie, ale może także poruszać inne momenty mszy. Najważniejszą rzeczą w tej części jest wydobycie konkretnych postaw, które postaramy się wprowadzać w nasze życie po powrocie. To spotkanie ma przede wszystkim charakter dzieleniowy. Należy „pomagać” uczestnikom w przypomnieniu ważnych fragmentów Eucharystii, a nawet szukać dziury w całym, prowokować do wypowiedzenia swojej opinii.


\section{Eucharystia jako zbawcza obecność Jezusa we wspólnocie wiernych}
\label{babice2006-jesien-knurow/spotkanie3:eucharystia-jako-zbawcza-obecnosc-jezusa-we-wspolnocie-wiernych}\begin{quote}

W pierwszy dzień Przaśników, kiedy ofiarowywano Paschę, zapytali Jezusa Jego uczniowie: «Gdzie chcesz, abyśmy poszli poczynić przygotowania, żebyś mógł spożyć Paschę?» I posłał dwóch spośród swoich uczniów z tym poleceniem: «Idźcie do miasta, a spotka was człowiek, niosący dzban wody. Idźcie za nim i tam, gdzie wejdzie, powiecie gospodarzowi: Nauczyciel pyta: gdzie jest dla Mnie izba, w której mógłbym spożyć Paschę z moimi uczniami? On wskaże wam na górze salę dużą, usłaną i gotową. Tam przygotujecie dla nas». Uczniowie wybrali się i przyszli do miasta, gdzie znaleźli, tak jak im powiedział, i przygotowali Paschę.

\begin{flushright}
---Mk 14, 12-16
\end{flushright}
\end{quote}

Podczas ostatniej wieczerzy wokół Jezusa zgromadziło się 12 apostołów. Oni to byli zalążkiem Kościoła, wspólnoty, jaką możemy obecnie przeżywać. Przed spożyciem wieczerzy Jezus nakazuje uczniom przygotować salę. Ma być „duża, usłana i gotowa”. Także i my powinniśmy przygotować się na spotkanie z Bogiem.
\begin{itemize}
\item {} 
Jak wygląda moje przygotowanie do Eucharystii? (Umawiając się na ważne spotkanie z kimś bliskim naszemu sercu, solidnie się przygotowujemy – wygląd, punktualność, przebieg rozmowy.)

\end{itemize}

Przeczytajmy:
\begin{quote}

Chleb, który łamiemy, czyż nie jest udziałem w Ciele Chrystusa? Ponieważ jeden jest chleb, przeto my, liczni, tworzymy jedno Ciało. Wszyscy bowiem bierzemy z tego samego chleba

\begin{flushright}
---1Kor 10, 16-17
\end{flushright}
\end{quote}
\begin{quote}

Czym w rzeczywistości jest chleb? Jest Ciałem Chrystusa. Kim stają się ci, którzy go przyjmują? Ciałem Chrystusa; ale nie wieloma ciałami, lecz jednym ciałem. Faktycznie, jak chleb jest jednością, choć składa się nań wiele ziaren, które choć się nie znają, w nim się znajdują, tak że ich różnorodność zanika w ich doskonałym zjednoczeniu — w ten sam sposób również my jesteśmy wzajemnie ze sobą zjednoczeni, a wszyscy razem z Chrystusem

\begin{flushright}
---Komentarz św. Jana Chryzostoma
\end{flushright}
\end{quote}

Poprzez Eucharystię, a zwłaszcza przyjęcie Jezusa w Komunii św. Sprawia, że stajemy się jedno z wszystkimi, którzy to Ciało przyjmują. Nasza różnorodność buduje wspólnotę.
\begin{itemize}
\item {} 
Jaka jest moja relacja, stosunek do pozostałych uczestników Mszy Św.?

\item {} 
Czy potrafię przyjąć i pogodzić się z róznorodnością ludzi uczestniczących w Eucharystii?

\item {} 
Co oznacza dla mnie, że ta róznorodność zanika w Jezusie? Jak rozumiem te słowa?
\begin{quote}

Jeśli więc przyniesiesz dar swój przed ołtarz i tam wspomnisz, że brat twój ma coś przeciw tobie, zostaw tam dar swój przez ołtarzem, a najpierw idź i pojednaj się z bratem swoim! Potem przyjdź i dar swój ofiaruj!

\begin{flushright}
---Mt 5, 23-24
\end{flushright}
\end{quote}

\end{itemize}

Jezus wzywa nas, abyśmy potrafili sobie wzajemnie wybaczać. Abyśmy mogli bez gniewu uczestniczyć w Ofierze Chrustusa. Żydzi zanim usiedli do wieczerzy z rodziną, przyjaciółmi obmywali sobie ręce. Oprócz względów higienicznych sprawiało to, że trzeba było rozluźnić rękę. Nie można już było mieć ręki zaciśnietej w pięść. Rozluźnienie ręki powodowało, że wszystkie złe emocje odchodziły, człowiek się uspokajał. Obecnie tą samą funkcję, a może jeszcze większą wyraża przekazanie „znaku pokoju”. Poprzez ten gest jednoczymy się z innymi ludźmi, przekazujemy sobie wzajemnie pokój, pokój chrystusowy – czyli uwolnienie od wszelkich zmartwień, trosk, złości, gniewu itp.
\begin{itemize}
\item {} 
Czy potrafię przebaczyć i pojednać się zanim wejdę do kościoła?

\item {} 
Jaką rolę w moim przeżywaniu Eucharystii pełni „znak pokoju”? Czy jest to tylko podanie ręki, bo tak wypada? Czy potrafię na mszy podać rękę tym, których skrzywdziłem, tym, którzy mnie skrzywdzili?

\end{itemize}


\section{Udział musi być czynny!}
\label{babice2006-jesien-knurow/spotkanie3:udzial-musi-byc-czynny}
Uczestnictwo we mszy musi być czynne. Nie można być tylko słuchaczem, kimś, kto przupatruje się z boku. Aby spotkanie mogło nastąpić, aby było owocne i korzystne dla obu stron, musi nastąpić czynny udział tych 2 stron.
\begin{itemize}
\item {} 
Jakie są moje odpowiedzi na wezwania kapłana?

\item {} 
Czy w ogóle odpowiadam, czy nie przysłuchuję się tylko co mówią inni? Jakie jest moje zaangażowanie w śpiew podczas liturgii?

\end{itemize}

Przeczytajmy:
\begin{quote}

Dar Chrystusa i Jego Ducha, który otrzymujemy w Komunii eucharystycznej, wypełnia z obfitością gorące pragnienia braterskiej jedności, jakie kryją się w sercach ludzkich, i jednocześnie wynosi doświadczenie braterstwa właściwego wspólnemu uczestnictwu w tym samym stole eucharystycznym na poziom daleko wyższy od zwyczajnego ludzkiego ucztowania. Dzięki zjednoczeniu z Ciałem Chrystusa Kościół coraz głębiej staje się w Chrystusie «jakby sakramentem, czyli znakiem i narzędziem wewnętrznego zjednoczenia z Bogiem i jedności całego rodzaju ludzkiego».

Jak pokazuje codzienne doświadczenie, początkom rozdziału między ludźmi, tak bardzo zakorzenionego w ludzkości z powodu grzechu, przeciwstawia się odradzająca jedność moc Ciała Chrystusa. Eucharystia, budując Kościół, właśnie dlatego tworzy komunię pomiędzy ludźmi.

\begin{flushright}
---Jan Paweł II -- Ecclesia de Eucharistia
\end{flushright}
\end{quote}

\emph{Warto tutaj zapytać uczesntników o uroczystości rodzinne (lub jakiekolwiek inne), które mają miejsce przy stole. Jak one wyglądają, jak zachowują się ludzie, którzy na nie przychodzą. Jak takie uroczystości są odbierane przez nas, przez uczestników (nie podejmować tematu alkoholu). Co sprawia, że ludzie zasiadają przy jednym stole, gromadzą się na uroczystości? Dlaczego to robią? Jaki jest cel tego spotkania?
Czy podobny cel ma zgromadzenie eucharystyczne? Co podczas niego nas przyciąga? Co sprawia, że gromadzimy się w jednym miejscu i czasie?}

Podczas Eucharystii stajemy się braćmi, jednoczymy się, nie sami z siebie, ale dzięki uczestnictwu w jednym Ciele i Krwi. Uczta Pańska nie jest zwykła ucztą. Od zawsze ludzie, aby być bliżej siebie, aby się jednoczyć organizują uczty. Najważniejsze wydarzenia w
życie każdego człowieka, radość, smutek przeżywamy właśnie podczas uczty, kolacji itp. Tam ludzie się spotykają, otwierają na siebie wzajemnie. Także te spotkania z osobą najbliższą odbywają się przy stole. Jednym z elementów doskonałego, romantycznego spędzenia czasu z ukochanym jest właśnie wspólna kolacja.

Eucharystia przenosi nas jednak na wyższy poziom. Tutaj jednoczy nas pokarm, który wspólnie spożywamy. To On pozwala nam dostrzegać jedność między nami, to On nas jednoczy. Należy zauważyć, że wszyscy ludzie gromadzą się dla tego Pokarmu, to on jest tym, co nas jednoczy. Uczta Pańska stanowi źródło wzajemnego zrozumienia,
poszanowania, źródło wspólnoty ludzi, odradza jedność między ludźmi, którą rozbił grzech. Eucharystia uczy nas miłości do naszych braci:
\begin{quote}

Eucharystia do tej miłości nas wychowuje w sposób najgłębszy, ukazuje bowiem jaką wartość w oczach Bożych ma każdy człowiek, nasz brat i siostra, skoro każdemu w taki sam sposób Chrystus daje siebie samego pod postaciami chleba i wina.

\begin{flushright}
---Jan Paweł II
\end{flushright}
\end{quote}


\section{Zastosowanie}
\label{babice2006-jesien-knurow/spotkanie3:zastosowanie}
Spróbujmy w tej części spotkania zastanowić się nad praktycznymi rzeczami wynikającymi z Eucharystii. Ponieważ Eucharystia nie jest oderwana od rzeczywistości, a ma nam pomagać w naszym codziennym życiu, jest obrazem idealnego życia we wspólnoscie z innymi ludźmi i Bogiem, dlatego zastanówmy się na podst. kazania ks. Blachnickiego, jakie elementy Eucharystii mogę i chcę poprawić i pogłębić w moim życiu (w kontekście tematu dnia). Przeanalizujmy wspólnie z uczestnikami tekst tego kazania. Spróbujmy wyakcentować te punkty liturgii, które możemy przenieść na grunt naszego życia codziennego.
\begin{quote}

W liturgii mamy nie tylko zjednoczyć się z Bogiem, ale zjednoczyć się między sobą i przeżyć w znaku zewnętrznym tę prawdę, że jesteśmy wspólnotą, że jesteśmy Kościołem, że jesteśmy ludem Bożym. I znów zależy to od nas, od naszego otwarcia się nie tylko ku Bogu, ale także ku braciom. {[}...{]} Kiedy przychodzimy na zgromadzenie liturgiczne to przychodzimy nie po to, żeby spełnić swój osobisty obowiązek religijny i nie po to, żeby osobiście się wzbogacić słowem Bożym czy łaską Bożą. Liturgia nie jest jakimś, można by użyć tego porównania, sklepem, gdzie każdy przychodzi, żeby zakupić dla siebie jakiś towar, który jest mu potrzebny. W liturgii nie ma tego, że ustawiamy się w ogonku i każdy bierze to, co jest jego i odchodzi zadowolony do domu. Liturgia w całości jest znakiem tego, czym jest Kościół, to znaczy znakiem wspólnoty. Idąc na mszę świętą, czy na inne zgromadzenie liturgiczne, idziemy po to, żeby łączyć się we wspólnotę, żeby dać wyraz temu, że każdy z nas jest członkiem wspólnoty, żeby przeżyć swoją jedność z całą wspólnotą Kościoła i żeby ją zewnętrznie wyrazić. W zgromadzeniu liturgicznym muszę poczuć się członkiem wspólnoty Kościoła, muszę wejść z całą świadomością w to zgromadzenie, w istotę chrześcijaństwa i Kościoła, to znaczy we wspólnotę z Bogiem i z braćmi.

Dlatego warunkiem wejścia w zgromadzenie liturgiczne jest pojednanie z braćmi, odrzucenie ze swojego serca tego wszystkiego, co nas oddziela. ``Jeżeli przynosisz dar swój do ołtarza - mówi Chrystus - a wspomnisz sobie po drodze, że brat twój ma coś przeciwko tobie, to idź najpierw pojednaj się z twoim bratem, a potem dopiero przyjdź i złóż Bogu ofiarę''. Otóż to słowo Chrystusa odnosi się przede wszystkim do naszego uczestnictwa w zgromadzeniu Eucharystycznym. Nikt z nas nie powinien wchodzić w to zgromadzenie, jeżeli jest skłócony z drugim człowiekiem, jeżeli ma w sercu do kogoś pretensje, jeżeli kogoś nienawidzi, jeżeli jeszcze komuś nie przebaczył, nie darował krzywdy.

Chrystus daje nam siebie i czeka na naszą odpowiedź wiary i miłości. Przyjmując komunię, czyli przyjmując Ciało Chrystusa, wyznaję przez to, że jestem w jedności z Chrystusem i ze wszystkimi, którzy przy tym samym stole ten sam Chleb pożywają. To jest znak jedności. Podobnie jak w naszym ludzkim życiu zasiadanie przy wspólnym stole jest znakiem przyjaźni, znakiem miłości. Ktoś, kto jest wrogo usposobiony, kto ma w sercu nienawiść, nie będzie proszony do stołu, nie zasiądzie przy wspólnym stole w gronie rodziny. Podobnie Eucharystia: spożywanie Ciała Pańskiego we wspólnocie ludu Bożego, sakrament - znak jedności. Jeżeli przystępuję do Stołu Pańskiego, przyjmuję Ciało Chrystusa, to tym samym wyrażam na zewnątrz to, że jestem w jedności z Chrystusem i ze wszystkimi, którzy są też w jedności z Chrystusem, którzy z tego samego Chleba pożywają. Jeżeli teraz przystępuje ktoś do tej zewnętrznej, sakramentalnej, komunii, a w sercu nie ma wiary ani miłości wobec Chrystusa,  nie  ma  w  swoim sercu  miłości wobec braci, nie  jest gotów wszystkim przebaczyć, nie jest ze wszystkimi pojednany, to spełnia wtedy zewnętrznie znak, który jest zakłamaniem, któremu nie odpowiadają wewnętrzne przeżycia i postawy

\begin{flushright}
---ks. Franciszek Blachnicki -- Fragmenty konferencji o liturgii skierowanej do uczestników oazy w Krościenku w dniu 28 VII 1973 r.
\end{flushright}
\end{quote}

Zapytajmy:
\begin{itemize}
\item {} 
Jak zmienia się nasza postawa po wyjściu z mszy św.? Na jak długo starcza nam tej siły do „znoszenia” innych ludzi we wspólnocie?

\item {} 
Czy Eucharystia ma jakieś punkty wspólne z naszym życiem codziennym? Czy nie jest ona może czymś zupełnie od niego oderwanym, czasem, w którym zachowujamy się zupełnie inaczej niż w rzeczywistości i nie staramy się przenieść tych wartośći do naszego żcyia?

\item {} 
W jaki sposób mogę lepiej przygotować się do mszy świętej? Czy stwarzam sobie możliwość takiego przygotowania?

\item {} 
Czy podczas mszy jestem otwarty na działanie Boga we mnie? Czy jestem ostwarty na innych uczestników liturgii?

\end{itemize}

Niech postanowieniem po tym spotkaniu będzie znalezienie jednej rzeczy na Eucharystii, którą będę starał się bardziej realizować po powrocie.(Np. niech to będzie znalezienie czasu na przygotowanie do mszy, głębsze spojrzenie na „znak pokoju”, próba zrozumienia innych uczestników liturgii, a nie tylko narzekanie na nich.) Natomiast podczas najbliższej Eucharystii podziękujmy Bogu za tę wspólnotę, która dokonuje się przez przyjęcie Ciała i Krwi Chrystusa. Podziękujmy Bogu za ludzi, za braci, którzy razem ze mną stanowią Ciało Jezusa Chrystusa.



\renewcommand{\indexname}{Indeks}
\printindex
\end{document}
